\section{Graph towers and spanning-tree valuations}\label{sec:gra}

\providecommand{\grah}{\mathsf{h}}
\providecommand{\graA}{\mathcal{A}}
\providecommand{\graP}{\mathbb{P}}
\providecommand{\gramu}{\boldsymbol{\mu}}
\providecommand{\graPi}{\Pi}
\providecommand{\gracut}{\vartheta}
\providecommand{\graunit}{\operatorname{unit}_p}
\providecommand{\grafiles}[1]{\begin{center}\footnotesize#1\end{center}}
\providecommand{\graJac}{\operatorname{Jac}}
\providecommand{\grahat}[1]{\widehat{#1}}

The results of this section are unconditional and use no hypothesis on zeros
of $L$-functions. They take from the source proof (\cref{src:sec:mechanism})
one local device: the exact valuation of $\zeta-1$ at a
root of unity $\zeta$ of prime-power order, together with the comparison of
the valuations of the terms of a polynomial in this uniformizer: a unique
term of least valuation cannot be cancelled. In a tower of graphs with
abelian deck groups this reads off the $p$-adic valuation of a Laplacian
eigenvalue at every character, and the valuations can be summed. For the
prime-to-$p$ part of the spanning-tree count, \cref{gra:sec:units} replaces
individual valuations by Galois stabilisers and the cyclotomic reciprocity
map, and treats towers with nonabelian $p$-adic analytic deck groups. The
other graph results, on quantum walks, signed square cycles and edge-polytope
zeros, are in \cref{sec:grw}; its Grover walks also give tower valuations of
the kind studied here (\cref{gra:thm:grover}).

\subsection{Voltage towers and conductor shells}\label{gra:sec:setting}
{\raggedright
\sources{grh_method_transfers_graph_towers_and_finite_rings; Effective_Graph_Towers; grh_method_transfers}
}

Fix a prime $p$ and a rank $d\ge1$, and normalise $v_p(p)=1$. A
\emph{voltage} is a vector attached to an oriented edge of a base graph; it
gives the displacement between the copies of its endpoints in a covering
graph. Let $G$ be a finite connected undirected multigraph with integral
voltages $a_e\in\Z^d$, one orientation per edge. At level $n\ge0$ the deck
group is $G_n=(\Z/p^n\Z)^d$, the vertices are the pairs $(x,g)$ with $x$ a
vertex of $G$ and $g\in G_n$, and an edge $e$ from $x$ to $y$ lifts to the
edges from $(x,g)$ to $(y,g+a_e)$. We assume throughout that every finite
lift is connected. Write $\kappa_n$ for the number of spanning trees at
level $n$ and
\[
 \nu_n=v_p(\kappa_n),\qquad \nu_0=v_p(\kappa(G)).
\]
The main case is a \emph{bouquet}: one vertex and $m$ loops with steps
$s_1,\dots,s_m\in\Z^d$. At level $n$ the vertex set is $G_n$, and for every
$x\in G_n$ and every listed step there is one undirected edge between $x$
and $x+s_i$. Repetitions count: parallel edges are retained, including both
edges if a step has order two, and a zero step gives loops, which contribute
nothing to the Laplacian. If the reductions of the steps span $\F_p^d$, they
generate every $G_n$, so all levels are connected; then $\nu_0=0$.

DuBose and Valli\`eres prove that $\nu_n$ eventually agrees with a
polynomial in $p^n$ and $n$ \cite[Theorem~6.2]{DuBoseVallieres2023Towers}.
The existence theorem does not specify the polynomial or the first level
where it holds; these are the subject of
\cref{gra:sec:effective,gra:sec:explicit,gra:sec:moments,gra:sec:extremal,gra:sec:onset}.

\paragraph{The shell identity.}
The \emph{Laurent Laplacian} of $(G,a)$ is the matrix over
$\Z[X_1^{\pm1},\dots,X_d^{\pm1}]$ indexed by the vertices of $G$ which, for
each edge $e$, adds $1$ to both endpoint diagonals and subtracts
$X^{a_e}$ and $X^{-a_e}$ in the two corresponding off-diagonal positions
(all four entries are at the same position for a loop). Let $F$ be its
determinant. For a bouquet,
\[
 F(X)=\sum_{i=1}^m\grah(X^{s_i}),\qquad
 \grah(U)=2-U-U^{-1}=-\frac{(U-1)^2}{U}.
\]
Character decomposition of the lifted Laplacian gives one block per
character of $G_n$, namely the Laurent Laplacian evaluated at the
corresponding tuple of $p^n$-th roots of unity; the trivial block is the
base Laplacian. A connected base has a connected lift at a given level
exactly when every nontrivial character block there is invertible, because
the multiplicity of the zero eigenvalue counts components. The matrix-tree
theorem applied to the base and to the lift cancels their common
vertex-count factor and leaves the deck-group order in the denominator.
Embedding the resulting algebraic identity into $\overline{\Q}_p$ and
grouping the characters by conductor yields
\begin{equation}\label{gra:eq:shells}
 \nu_n=\nu_0-dn+\sum_{h=1}^n S_h,\qquad
 S_h=\sum_{\xi}v_p(F(\xi)),
\end{equation}
where $\xi$ runs over the tuples of $p$-power roots of unity of joint order
exactly $p^h$. The same argument works for graphs with leaves and for
cycles.

Fix a primitive $p^h$-th root of unity $\zeta=\zeta_h$ and write
$e=e_h=\varphi(p^h)=(p-1)p^{h-1}$ and $\pi=\zeta-1$. A point of joint
conductor $p^h$ is $\zeta^u=(\zeta^{u_1},\dots,\zeta^{u_d})$ with
$u\in(\Z/p^h\Z)^d$ \emph{primitive}, that is, with at least one unit entry.
The simultaneous action of the units modulo $p^h$ on such tuples is free;
its orbits have size $e$ and consist of cyclotomic conjugates, so the
valuation is constant on an orbit. Hence
\begin{equation}\label{gra:eq:orbits}
 S_h=\sum_{[u]}\ord_h F(\zeta^u),\qquad \ord_h(x)=e_h\,v_p(x),
\end{equation}
the sum being over orbits. We call $\ord_h$ the \emph{normalised order} at
conductor $p^h$. This orbit decomposition is already used by DuBose and
Valli\`eres; the point here is to evaluate the individual orders.

\begin{lemma}[Comparison rules]\label[lemma]{gra:lem:rules}
The polynomial $\Phi_{p^h}(1+T)$ is Eisenstein, so $v_p(\pi)=1/e$. Let
$f\in\Z_p[[T]]$.
\begin{enumerate}[label=\textup{(\roman*)}]
\item If the first unit coefficient of $f$ is at degree $k<e$, then
$\ord_h f(\pi)=k$.
\item The same conclusion holds for $k<e+2$ if the constant and linear
coefficients of $f$ vanish identically.
\item For arbitrary $f\in\Z_p[T]$, reduce $f$ modulo $\Phi_{p^h}(1+T)$ to
$\sum_{j<e}c_jT^j$. If this is nonzero, then
$\ord_h f(\pi)=\min_{c_j\ne0}\{e\,v_p(c_j)+j\}$.
\end{enumerate}
\end{lemma}

\begin{proof}
Compare valuations term by term. In (i) the lower coefficients are
divisible by $p$ and contribute normalised order at least $e>k$, and higher
powers have order greater than $k$. In (ii) a lower nonzero term has degree
at least two and coefficient divisible by $p$, hence order at least $e+2$.
In (iii) the integers $e\,v_p(c_j)+j$ lie in distinct residue classes
modulo $e$, so the minimum is attained once.
\end{proof}

The square-torus symbol $\grah(X)+\grah(Y)$ recurs below, and its valuation
is immediate when its residual quadratic form is anisotropic.

\begin{lemma}[Anisotropic sums of two squares]\label[lemma]{gra:eq:aniso}
Let $p$ be a prime such that $-1$ is not a square modulo $p$, and let
$(U,V)$ be a pair of $p$-power roots of unity of joint conductor $p^j>1$.
Then
\[
 v_p\bigl(\grah(U)+\grah(V)\bigr)=2/e_j\le1 .
\]
\end{lemma}

\begin{proof}
Write $U=\zeta^a$, $V=\zeta^b$ in a common primitive coordinate $\zeta$ of
order $p^j$, with $a,b$ not both divisible by $p$. Since
$\grah(\zeta^a)=-(\zeta^a-1)^2\zeta^{-a}$ and
$(\zeta^a-1)/(\zeta-1)\equiv a$ modulo $\zeta-1$, the reduction of
$(\grah(U)+\grah(V))/(\zeta-1)^2$ modulo $\zeta-1$ is $-(a^2+b^2)$. Since
$-1$ is a nonsquare modulo $p$, this is nonzero, and the valuation is
$2v_p(\zeta-1)=2/e_j$. Finally $p$ is odd, so $e_j\ge2$.
\end{proof}

For a general collected Laurent polynomial the orbit order is computed as
follows. Write
\begin{equation}\label{gra:eq:laurent}
 F(X)=p^g\sum_{i=1}^m c_iX^{a_i},\qquad a_i\in\Z^d\text{ distinct},\quad
 c_i\in\Z\setminus\{0\},\quad \min_iv_p(c_i)=0 .
\end{equation}
Put $Q=p^{h-1}$, so that $e=(p-1)Q$. For each monomial choose
$0\le n_i<pQ$ congruent to $a_i\cdot u$ modulo $pQ$, and write
$n_i=r_i+Qb_i$ with $0\le r_i<Q$. Set
$A_{rb}=\sum_{i:(r_i,b_i)=(r,b)}c_i$. The unique representative of the
normalised evaluation $p^{-g}F(\zeta^u)$ in the power basis
$1,\zeta,\dots,\zeta^{e-1}$ is $R(\zeta)$, where
\[
 R(Y)=\sum_{r=0}^{Q-1}\sum_{b=0}^{p-2}C_{rb}Y^{r+Qb},\qquad
 C_{rb}=A_{rb}-A_{r,p-1},
\]
because $\Phi_{p^h}(Y)=1+Y^Q+\dots+Y^{(p-1)Q}$ eliminates the last block of
exponents. In particular $R=0$ detects exact vanishing.

\begin{lemma}[Integral orbit valuation]\label[lemma]{gra:lem:orbit}
If $R\ne0$, let $\mu=\min_{C_{rb}\ne0}v_p(C_{rb})$, and let $j$ be the
first index whose coefficient in $p^{-\mu}R(1+T)$ is nonzero modulo $p$.
Then $0\le j<e$ and
\[
 \ord_hF(\zeta^u)=e(g+\mu)+j .
\]
\end{lemma}

\begin{proof}
Changing variables from $Y$ to $1+T$ preserves nonzero polynomials over
$\F_p$, so $j$ exists. In $p^{-\mu}R(1+\pi)$ the term of degree $j$ has
valuation $j/e<1$. Every earlier nonzero coefficient is divisible by $p$
and contributes valuation at least $1$, and each later term has valuation
strictly greater than $j/e$. The unique term of least valuation gives the
assertion, also when $p=2$ and $h=1$.
\end{proof}

The integer content is uniformly bounded: the sums defining the two terms
of $C_{rb}$ use disjoint subsets of the coefficients, so
$p^\mu\le\sum_i\abs{c_i}$ when $R\ne0$. The content must be computed after
collecting terms. For $p=2$, $h=2$ and the square-grid symbol at $u=(1,1)$
the exponent polynomial $4-2X-2X^3$ reduces to $4$ modulo $X^2+1$; thus
$g=0$, $\mu=2$, $j=0$ and the orbit contributes $4$. Computing the content
before collecting the terms would give the wrong answer.

\subsection{The eventual polynomial and its first valid level}\label{gra:sec:effective}
\sources{Effective_Graph_Towers}

\begin{theorem}[Effective polynomial and first valid level; \stU]\label{gra:thm:effective}
Let $(G,a)$ be a finite connected multigraph with integral voltages in
$\Z^d$ all of whose lifts with deck groups $(\Z/p^n\Z)^d$ are connected;
for instance a bouquet whose steps span $\F_p^d$ after reduction. From $p$
and the voltages one can compute the unique rational polynomial $P(X,Y)$ of
total degree at most $d$ and of degree at most one in $Y$, and the least
$N\ge0$, such that $\nu_n=P(p^n,n)$ for every $n\ge N$. The algorithm also
returns the exact rational generating function of all $\nu_n$ and the
finite polynomial of exceptions. If levels start at one, the least start is
$\max(1,N)$. The same holds for rational voltages with denominators prime
to $p$.
\end{theorem}

The proof treats a finite Laurent polynomial \eqref{gra:eq:laurent} before
specialising to graph Laplacians, and evaluates all conductors at once: it
gives a finite representation of the first surviving coefficient of
\cref{gra:lem:orbit} and shows that its output yields the entire generating
function. It occupies the rest of this subsection.

\paragraph{Reading the orbits digit by digit.}
Normalise each orbit representative $u$ by making its first unit coordinate
equal to $1$ modulo $p^h$. In the chart with pivot $i_0$, lower-indexed
coordinates have units digit zero; the units digit of the pivot is one and
its other digits are zero; every other digit is free. Reading the
coordinate vectors from the least significant end therefore enumerates each
orbit exactly once. At every stage there are two operations on a new digit:
treating it as the top digit emits a valuation at that conductor, and
treating it as an ordinary digit prepares the next stage.

After $k$ ordinary digits let $Q=p^k$, let $w\in\{0,\dots,Q-1\}^d$ be the
prefix, and decompose
\[
 a_i\cdot w=r_i+Qs_i,\qquad 0\le r_i<Q .
\]
Only the carries $s_i$ and a small amount of information about the $r_i$
are needed. For $0\le j<Q$ define a row vector over $\F_p$ by
\[
 B_j=\Bigl(\tbinom{r_1}{j},\dots,\tbinom{r_m}{j}\Bigr)\pmod p .
\]
Process these rows in order of increasing $j$, retaining a row precisely
when it is outside the span of the earlier retained rows. Store the ordered
list of retained vectors, and keep their integer indices in separate
registers. There are at most $m$ such registers.

\begin{lemma}[What the retained rows determine]\label[lemma]{gra:lem:lucas}
Two columns agree on the retained rows if and only if their residues $r_i$
agree. For any linear functional on the rows, its first nonzero value, if
one exists, occurs on a retained row. The retained list after one more
ordinary digit is computable from the previous list and the carries.
\end{lemma}

\begin{proof}
The full column associated to $r$ is the coefficient vector of
$(1+T)^r\in\F_p[T]$. Distinct $r<Q$ give different degrees, and agreement
on a basis of the row space implies agreement on all rows. A discarded row
lies in the span of earlier rows, so it cannot be the first detected by a
linear functional.

For a new coordinate digit vector $v\in\{0,\dots,p-1\}^d$, perform
Euclidean division $s_i+a_i\cdot v=b_i+ps_i'$ with $0\le b_i<p$. The new
residue is $r_i+Qb_i$ and the new carry is $s_i'$. Over $\F_p$ one has
$(1+T)^Q=1+T^Q$, hence
\[
 \binom{r_i+Qb_i}{j+Qt}\equiv\binom{r_i}{j}\binom{b_i}{t}\pmod p
 \qquad(0\le j<Q,\ 0\le t<p).
\]
For each $t=0,\dots,p-1$ in order, multiply every old retained row
coordinatewise by $(\binom{b_i}{t})_i$, and feed these candidates, in old
row order within the block, to the same greedy independence test. The
omitted old rows are combinations of earlier rows in their own block after
this multiplication, so inserting them into the scan would leave all
decisions unchanged. A retained candidate coming from an old index $j$ has
new index $j+Qt$.
\end{proof}

\paragraph{Reading the top digit.}
Use the carries to compute $b_i$ as above, and partition the columns
according to their entries on the stored basis. \Cref{gra:lem:lucas}
identifies these groups with the occupied values of $r$. Compute the
integers $C_{rb}$ in each group, test for vanishing, and then compute
$\mu$ as in \cref{gra:lem:orbit}. Division by $p^\mu$ takes place
\emph{after} the integer sums and subtractions within each group. In top
block $t=0,\dots,p-2$ put
\[
 w_r(t)=\sum_{b=0}^{p-2}(C_{rb}/p^\mu)\binom bt\pmod p .
\]
At index $j_0+Qt$, the coefficient to be tested modulo $p$ is
$\sum_rw_r(t)\binom r{j_0}$. This is a linear functional on the old rows,
using one coordinate per group. Testing the stored rows in order finds its
earliest nonzero value; stop at the first successful block. If the
successful old index is $j_0$, the orbit contribution is
\[
 \bigl((p-1)(g+\mu)+t\bigr)Q+j_0 .
\]
All choices except the numerical values of $Q$ and $j_0$ depend only on
the finite state and this last digit.

\paragraph{A reduction of the state space.}
After processing $k$ ordinary digits, test the normalised coefficient
functional $B_j\mapsto\sum_ic_i\binom{r_i}{j}$. Suppose its first nonzero
value is at $j_*<p^k$. For every subsequent conductor the ordinary exponent
polynomial at $1+T$, reduced modulo $(p,T^{p^k})$, equals
$\sum_ic_i(1+T)^{r_i}$ in this quotient, since further exponent digits
multiply these terms by powers of $1+T^{p^k}$. Final cyclotomic reduction
preserves the same truncation, because
$\Phi_{p^h}(1+T)\equiv T^{\varphi(p^h)}\bmod p$ and $\varphi(p^h)\ge p^k$
for $h>k$. Consequently every continuation has $\mu=0$ and $j=j_*$,
including arbitrary choices for its final digit. Such a prefix can be
replaced by a state storing only $Q$ and $j_*$, with the same chart; future
ordinary digits replace $Q$ by $pQ$ and leave $j_*$ fixed.

\paragraph{Termination and weights.}
The control state consists of the pivot, a first-digit flag, the carries
and the ordered list of row vectors; the register values are stored
separately. If $H_i=\sum_\ell\abs{a_{i\ell}}$, then
$\abs{s_i}=\abs{\floor{a_i\cdot w/Q}}\le H_i$ for every prefix. An explicit
upper bound for the number of control states, including the reduced
states, is
\begin{equation}\label{gra:eq:states}
 N_{\mathrm{st}}=d+2d\Bigl(\prod_{i=1}^m(2H_i+1)\Bigr)\sum_{s=0}^mp^{ms}.
\end{equation}
A breadth-first search enumerates the reachable states and tests every
allowed final digit. It therefore also decides whether a nontrivial
$p$-power torsion zero of $F$ exists. One can reject such inputs or, for a
sum taken only over nonzero evaluations, give a zero evaluation weight
zero while continuing the search; rejection is appropriate for an ordinary
product.

Each ordinary state has registers $(Q,j_1,\dots,j_s)$ with $s\le m$. Each
ordinary digit induces an integer linear map between registers by $Q'=pQ$,
$j_a'=j_b+tQ$. Summing the maps for digits with the same target gives the
corresponding block of a matrix $M$. In each state the sum of the output
forms over its legal top digits gives the corresponding block of a row
vector $f$. Initially each chart has zero carries, the single row
$(1,\dots,1)$ and registers $(1,0)$. Let $u_0$ be the concatenated initial
vector, zero outside these initial states. Matrix iteration adds the
registers of all prefixes at each state. By induction on the prefix length
and the unique chart representation, the shell sum at conductor $p^h$ is
exactly
\begin{equation}\label{gra:eq:transfer}
 S_h=fM^{h-1}u_0 .
\end{equation}
The matrix dimension $L$ is the sum of the register counts of the
reachable states; in particular $L\le(m+1)N_{\mathrm{st}}$.

\paragraph{The generating function.}
Let $F$ be the determinant of the Laurent Laplacian of $(G,a)$. With $z$ a
formal variable, summing \eqref{gra:eq:shells} and using
\eqref{gra:eq:transfer} yields
\begin{equation}\label{gra:eq:resolvent}
 \sum_{n\ge0}\nu_nz^n=\frac{\nu_0}{1-z}-\frac{dz}{(1-z)^2}
 +\frac z{1-z}\,f(I-zM)^{-1}u_0 .
\end{equation}
This identity computes the entire sequence. If
$\chi_M(x)=x^L+c_1x^{L-1}+\dots+c_L$, then
$\det(I-zM)=1+c_1z+\dots+c_Lz^L$. The adjugate shows that the numerator of
$f(I-zM)^{-1}u_0$ over this determinant has degree at most $L-1$.
Multiplying that denominator by the first $L$ shell coefficients of the
series and keeping degrees below $L$ gives its exact numerator, including
when $M$ is singular. The base cofactor and the Laurent determinant are
explicit finite input calculations.

\begin{proof}[Proof of \cref{gra:thm:effective}]
Put
\[
 D(z)=(1-p^dz)\prod_{i=0}^{d-1}(1-p^iz)^2 .
\]
An eventual law $\nu_n=P(p^n,n)$ with $\deg P\le d$ and $\deg_YP\le1$
implies that the reduced denominator of \eqref{gra:eq:resolvent} divides
$D(z)$. Conversely, whenever it divides $D(z)$, polynomial division and
partial fractions give
\[
 \sum_{n\ge0}\nu_nz^n=E(z)+\sum_{i=0}^d\Bigl(\frac{A_i}{1-p^iz}
 +\frac{B_i}{(1-p^iz)^2}\Bigr),\qquad B_d=0,
\]
with a polynomial $E$. Thus $P(X,Y)=\sum_i[A_i+(Y+1)B_i]X^i$ and
$\nu_n-P(p^n,n)=[z^n]E$. If $E=0$ the least onset is zero; otherwise it is
$1+\deg E$. Uniqueness of $P$ follows since the largest exponential base in
a nonzero sum $\sum_i(a_i+b_in)p^{in}$ dominates its lower bases on a tail.

DuBose and Valli\`eres guarantee the denominator shape for base graphs with
no leaves and nonzero Euler characteristic
\cite[Theorem~6.2]{DuBoseVallieres2023Towers}. To pass to an arbitrary
connected base admitting such a connected tower, repeatedly remove pendant
edges: their lifts are pendant edges and belong to every spanning tree. If
the remaining core is a cycle, its fundamental group is cyclic, so a
connected tower with deck groups as above has $d=1$; a cycle of length $b$
then lifts to a cycle of length $bp^n$, with $\nu_n=v_p(b)+n$ (including a
single loop or a double edge). A tree admits no connected nontrivial finite
cover. These cases exhaust the missing base graphs. Computed divisibility
also certifies the shape directly in each example, and for a general
Laurent sum the characteristic-polynomial procedure produces its exact
rational series independently of the existence theorem.

Finally, multiplying every voltage in $\Q\cap\Z_p$ by a common denominator
prime to $p$ induces compatible deck-group automorphisms at every level, so
the finite graphs are isomorphic to those with the cleared integral
voltages.
\end{proof}

\begin{remark}[A bound for the onset]\label[remark]{gra:rem:onsetbound}
Formula \eqref{gra:eq:resolvent} has numerator and denominator degrees at
most $L+2$. Consequently $\deg E\le L+2$ if $E\ne0$, and the least onset is
at most $L+3$, where $L\le(m+1)N_{\mathrm{st}}$ with $N_{\mathrm{st}}$ as in
\eqref{gra:eq:states}. The computation itself requires only the first $L$
shell sums after the reachable matrix has been built. This bound is very
far from the truth in rank two, where the onset is at most logarithmic in
the size of the voltages (\cref{gra:cor:logbound}); no polynomial bound for
the running time is claimed.
\end{remark}

\subsection{Explicit towers: the examples of DuBose--Valli\`eres and of Kataoka}\label{gra:sec:explicit}
\sources{grh_method_transfers_graph_towers_and_finite_rings; Effective_Graph_Towers}

Section~7 of \cite{DuBoseVallieres2023Towers} fits candidate polynomials to
finite data for five two-dimensional bouquet towers. Kataoka's preprint
\cite[\S6.3, Tables~6--7]{Kataoka2026Iwasawa} treats dyadic towers with
prescribed leading coefficient $\lambda$; its Table~7 lists data through
level six and suggests a full formula for each $0\le\lambda\le8$, of which
the coefficient of $n2^n$ is verified there and the remaining coefficients
are left unconfirmed. The following two theorems determine all fourteen
formulas and their first valid levels. Throughout, $e_1=(1,0)$,
$e_2=(0,1)$, and for $u=(a,b)$ primitive modulo $p^h$ we write
$X=\zeta^a$, $Y=\zeta^b$ and $F_u(\zeta)=F(\zeta^a,\zeta^b)$.

\begin{theorem}[The five towers of DuBose and Valli\`eres; \stU]\label{gra:thm:five}
The five towers of \cite[\S7]{DuBoseVallieres2023Towers} have the following
valuations. The column $N$ gives the least level, with level zero included,
from which the formula holds.
\begin{center}\small\upshape
\begin{tabular}{cclll}
\toprule
Example & $p$ & steps & $\nu_n$ & $N$\\
\midrule
1 & 2 & $e_1,e_2$ & $(2n+4)2^n-6n-1$ & 1\\
2 & 2 & $e_1,e_1,e_2,e_2$ & $4^n+(2n+4)2^n-6n-2$ & 1\\
3 & 2 & $(1,5),(0,3),(1,2),(0,1)$ & $(n+33/4)2^n-4n-1$ & 4\\
4 & 3 & $e_1,e_2$ & $4\cdot3^n-2n-4$ & 0\\
5 & 3 & $(1,0),(2,3),(1,1)$ & $20\cdot3^{n-1}-2n-8$ & 1\\
\bottomrule
\end{tabular}
\end{center}
In Example~3 the values at levels $0,1,2,3$ are $0,5,19,65$.
\end{theorem}

The values of $N$ follow from the formulas for $n\ge1$ (respectively
$n\ge4$) proved below: at $n=0$ the expressions in rows 1, 2 and 5 do not
vanish while that of row 4 does, and in row 3 the polynomial predicts $77$
at level three. The algorithm of \cref{gra:thm:effective}, supplied only
with the prime and the steps, returns the same five polynomials and the
same least levels; its integer linear representations have dimensions
$25,25,32,8,16$, and in Example~3 the finite correction is
$E(z)=-29/4-(17/2)z-13z^2-12z^3$
(\artifact{Effective_Graph_Towers/dv_examples.json}).

\begin{proof}[Proof for Examples 1 and 2]
Here $p=2$ and $F=4-X-X^{-1}-Y-Y^{-1}$. The identity
\[
 F=4-(X+Y)\bigl(1+(XY)^{-1}\bigr)
\]
makes the exceptional characters explicit. Let $h\ge3$ and $e=2^{h-1}$.
For a mixed-parity primitive pair, both $a-b$ and $a+b$ are odd, so each
factor in the product has valuation $1/e$, giving normalised order $2$.
There are $2e$ mixed-parity orbits, contributing $4e$.

For the remaining orbits normalise $a=1$, so $b$ runs through the $e$ odd
residues modulo $2^h$. The normalised orders of the factors are those of
$1+\zeta^{1-b}$ and $1+\zeta^{1+b}$. For an exponent $k$ put $r=k+e$ modulo
$2^h$. If $r\ne0$, the order of $1+\zeta^k$ is $2^{v_2(r)}$; if $r=0$ the
factor vanishes. For either sign, $1\pm b+e$ runs exactly once through the
even residues. Of a given pair of exponents, one is $2$ modulo $4$, so its
order is $2$; the other, when finite, has order at most $e$. Thus a finite
product has normalised order at most $e+2<2e$, and subtracting from $4$
preserves that order. At a vanishing factor the order of $F$ is $2e$
instead.

For either sign, the sum over the nonzero even residues is
$\sum_{i=1}^{h-1}(e/2^i)2^i=(h-1)e$. At each of the two exceptional orbits
the other factor has already contributed $2$ to this sum, so each requires
a correction of $2e-2$. The odd--odd contribution is $2(h-1)e+2(2e-2)$, and
consequently
\[
 S_h=(h+3)2^h-4\qquad(h\ge3).
\]
At $h=1$ the eigenvalues are $4,4,8$, giving $S_1=7$. At $h=2$ the four
mixed orbits each contribute $2$ and the two odd--odd orbits each
contribute $4$, giving $S_2=16$. Equation \eqref{gra:eq:shells} gives row 1
by summation. Doubling each edge multiplies the number of spanning trees by
$2^{4^n-1}$, since a tree uses $4^n-1$ edges with two independent parallel
choices each. This proves row 2.
\end{proof}

\begin{proof}[Proof for Example 3]
Write $F=8-\sum_{(u,v)}(X^uY^v+X^{-u}Y^{-v})$, the sum being over the four
steps. The two vectors $(1,2),(0,1)$ generate the group, so the only zero
of $F$ among $2$-power torsion points is the trivial character, and
$\nu_h-\nu_{h-1}=S_h-2$. The decisive local computation is the following
identity for the reduction $\overline F$ over $\F_2$:
\[
 XY^5\,\overline F=(Y+1)\bigl[(Y^2+Y+1)(X^2Y^7+1)+XY^2(Y+1)^5\bigr].
\]
Substitute $X=(1+z)^a$, $Y=(1+z)^b$, using nonnegative representatives of
$a,b$. For $b\ne0$ put $B=2^{v_2(b)}$ and $C=2^{v_2(2a+7b)}$. In
characteristic two, $(1+z)^j-1$ has order $2^{v_2(j)}$. Since $C$ is a
power of two, $C\ne5B$, and the resulting order of $\overline F$ is exactly
$D=B+\min(C,5B)$. Whenever $D<e$ this gives $\ord_hF(\zeta^a,\zeta^b)=D$ by
\cref{gra:lem:rules}(i), applied to the polynomial $XY^5F$ to avoid
negative exponents.

For $h\ge5$ and $e=2^{h-1}$, the following disjoint classes exhaust the
conductor shell.
\begin{center}\small\upshape
\begin{tabular}{lll}
\toprule
Class & normalised order & contribution to $S_h$\\
\midrule
$b$ odd & $2$ & $4e$\\
$b\equiv2\pmod4$ & $6,10,12$ & $17e/4$\\
$v_2(b)=k$, $2\le k\le h-2$ & $2^k+2$ & $(h-3)e+e-4$ (all $k$)\\
$b=0$ & $e+2$ & $e+2$\\
$b=e$ & $2e$ & $2e$\\
\bottomrule
\end{tabular}
\end{center}
Indeed, for fixed $b\equiv2\pmod4$, as $a$ runs through odd residues, $C$
is $4$, $8$, or at least $16$ in proportions $1/2,1/4,1/4$. For $k\ge2$ we
have $C=2$ and $e/2^k$ eligible values of $b$ at depth $k$. All three
generic rows have $D<e$. For the two exceptional rows, the integer
identities $F(X,1)=2(2-X-X^{-1})$ and $F(X,-1)=12$ give the listed orders
directly. Summation yields
\[
 S_h=(h+37/4)2^{h-1}-2\qquad(h\ge5).
\]
The small-shell totals $S_1,S_2,S_3,S_4$ are $7,16,48,116$. For an exact
short verification use one representative per Galois orbit, $(a,1)$ with
arbitrary $a$ and $(1,b)$ with $b$ even, and apply
\cref{gra:lem:rules}(iii). At $h=1,2,3,4$ the multisets of orbit
contributions, with superscripts for multiplicities, are respectively
\[
 \{2^{[2]},3\},\quad\{2^{[4]},4^{[2]}\},\quad
 \{2^{[8]},6^{[2]},8,12\},\quad\{2^{[16]},6^{[4]},10^{[2]},16,24\}.
\]
Thus $(\nu_1,\nu_2,\nu_3,\nu_4)=(5,19,65,179)$, and summing the shell
recurrence proves the formula for all $n\ge4$.
\end{proof}

\begin{proof}[Proof for Examples 4 and 5]
Row 4 is the case $p=3$, $m=2<p$ of \cref{gra:thm:cutoff} below: the
coefficient of $T^2$ in $L_a$ is $-(a_1^2+a_2^2)$, nonzero modulo three as in
\cref{gra:eq:aniso}, so $j(a)=2$ for all eight nonzero $a$, $C=8$ and
$\nu_n=4\cdot3^n-4-2n$ for $n\ge1$; at $n=0$ both sides vanish.

For row 5 the projected steps are $a$, $2a+3b$, $a+b$. At $h=1$ the
eigenvalue is three times the number of nonzero projected steps modulo
three. The four orbits have normalised orders $4,4,2,2$, so $S_1=12$. For
$h\ge2$, classify the primitive direction modulo nine. The following table
gives the first nonzero degree $k$ of $F_u(1+T)$ modulo $(3,T^9)$.
\begin{center}\small\upshape
\begin{tabular}{c|rrrrrrrrr|ccc}
$u$ & $(1,0)$ & $(1,1)$ & $(1,2)$ & $(1,3)$ & $(1,4)$ & $(1,5)$ &
$(1,6)$ & $(1,7)$ & $(1,8)$ & $(0,1)$ & $(3,1)$ & $(6,1)$\\
\hline
$k$&4&6&2&6&4&2&4&4&2&2&2&2
\end{tabular}
\end{center}
For an explicit verification put $J=\{a,2a+3b,a+b\}$ as a multiset. For
$1\le j\le6$ the coefficient is the reduction modulo three of
$-\sum_{k'\in J}\bigl(\binom{[k']_9}{j}+\binom{[-k']_9}{j}\bigr)$, where
the representatives lie in $[0,8]$; the constant is zero before reduction.
The table persists at every higher lift since $(1+T)^9\equiv1$ modulo
$(3,T^9)$. Normalising a unit multiple amounts to a change of variable with
invertible linear term. In the original Laurent expansion both the
constant and the linear terms vanish exactly; thus $k\le6<e+2$ makes each
displayed $k$ the normalised order at every $h\ge2$. Every displayed orbit
has $3^{h-2}$ orbit lifts, and the sum of the table is $40$, so
$S_h=40\cdot3^{h-2}$. Summing from $S_1=12$ proves row 5.
\end{proof}

In the next theorem $X,Y$ denote voltage monomials for the two basis
vectors, so $X^aY^b$ corresponds to the step $(a,b)$, and a repeated entry
means parallel generators.

\begin{theorem}[Kataoka's nine rows; \stU]\label{gra:thm:kataoka}
Let $p=2$ and $N=2^n$. For the nine step lists of
\cite[Table~7]{Kataoka2026Iwasawa} below, $\nu_n$ equals the given
expression for every $n$ from the displayed first level on.
\begin{center}\small\upshape
\begin{tabular}{cllc}
\toprule
$\lambda$ & generators & $\nu_n$ & first $n$\\
\midrule
0 & $X^3Y^2,X,Y$ & $9N-2n-3$ & 3\\
1 & $X^2Y^2,X,Y$ & $nN+(15/2)N-4n-7$ & 3\\
2 & $X,Y,Y$ & $2nN+2n+1$ & 2\\
3 & $XY,X,Y$ & $3nN+9N-8n-9$ & 1\\
4 & $X^2,X,X,Y,Y$ & $4nN-4N+6n+8$ & 2\\
5 & $X^2Y^2,X^2Y,XY^2,X,Y$ & $5nN+5N-16n+22$ & 3\\
6 & $X^3,X,Y,Y$ & $6nN-10N+14n+14$ & 3\\
7 & $X^3Y^3,X^3Y^2,X^2Y^3,X^2Y,XY^2,XY$ & $7nN+4N-20n+14$ & 3\\
8 & $X^4,X,X,Y,Y$ & $8nN-16N+14n+6$ & 2\\
\bottomrule
\end{tabular}
\end{center}
\end{theorem}

In every row the lifts are connected and the only $2$-power torsion zero of
the Laplacian symbol $F$ is the trivial character, so
\eqref{gra:eq:shells} reads $\nu_0=0$, $\nu_h=\nu_{h-1}+S_h-2$. In each
case the proof gives a closed formula for $S_h$ for large $h$ and the first
few shell sums by the exact reduction of \cref{gra:lem:rules}(iii); the
latter finite computations are recorded in
\grafiles{\artifact{grh_method_transfers_graph_towers_and_finite_rings/verification/kataoka_results.txt}\\
\artifact{grh_method_transfers_graph_towers_and_finite_rings/verification/odd_rows_results.txt}}
Write $e=2^{h-1}$.

\begin{proof}[Proof for $\lambda=0$]
For $u=(a,b)$ the projected steps are $3a+2b,a,b$. Reduce the expansion of
$F_u(1+T)$ modulo $(2,T^8)$, exactly as in the ternary table above. For
$(1,b)$, $0\le b\le7$, the orders are $6,2,4,2,6,2,4,2$ in sequence; for
$(a,1)$, $a=0,2,4,6$, they are all $2$. These orders are stable on lifting
each class modulo eight, and they are strictly smaller than $e$ for
$h\ge4$. The sum of the twelve entries is $36$, so $S_h=36\cdot2^{h-3}$.
The first three exact shell sums $7,18,44$ give $\nu_3=63$, and the
recurrence yields $\nu_n=9\cdot2^n-2n-3$ for $n\ge3$.
\end{proof}

\begin{proof}[Proof for $\lambda=1$]
The three steps are $(2,2),(1,0),(0,1)$. Put $Z=XY$. The eigenvalue has the
exact factorisation
\[
 F=4-(Z+1)H,\qquad H=(X+Y)/Z+(Z-1)^2(Z+1)/Z^2,
\]
since the third pair of translated terms adds $-(Z-Z^{-1})^2$ to the
square-torus eigenvalue. Work first with $h\ge4$. For mixed parities the
square-torus eigenvalue has order $2$ and $\ord_h((Z-Z^{-1})^2)=4$, so the
$2e$ mixed orbits contribute $4e$.

For odd $a,b$ take $a=1$. If $b\equiv-1\pmod4$, the first term of $H$ has
order $2$, since $\ord_h(X+Y)=\ord_h(1-\zeta^{b-1+e})=2$; the second has
order at least $12$ or vanishes. Thus $\ord_h(H)=2$. Over these $e/2$
orbits, the exponents $1+b+e$ range over all multiples of four modulo
$2^h$. For a nonzero such exponent $k$, the factor $Z+1$ has order
$2^{v_2(k)}\le e$. The finite orders sum to
$\sum_{j=2}^{h-1}(e/2^j)2^j=(h-2)e$. One orbit has $Z=-1$, giving
$\ord_h(F)=2e$; on the others $\ord_h(F)=\ord_h(Z+1)+2<2e$. These orbits
together contribute $(h+1)e-2$.

For $b\equiv1\pmod4$ both $Z-1$ and $Z+1$ have order $2$, so the second
term of $H$ has order $6$. The first has order $4$ for $b\equiv5\pmod8$ and
at least $8$ (including infinity) for $b\equiv1\pmod8$, since $8\mid e$.
The two contributions to $H$ have distinct orders in both cases. The
resulting orders of $F$ are $6$ and $8$, each occurring $e/4$ times; their
sum is $7e/2$. Hence
\[
 S_h=(h+17/2)2^{h-1}-2\qquad(h\ge4).
\]
Direct cyclotomic reduction at $h=1,2,3$ gives the shell sums $7,16,48$,
so $(\nu_1,\nu_2,\nu_3)=(5,19,65)$, and the recurrence gives
$\nu_n=(n+15/2)2^n-4n-7$ for $n\ge3$.
\end{proof}

\begin{proof}[Proof for $\lambda\in\{2,4,6,8\}$ and for $\lambda=3$]
For $\lambda=2\ell$, $1\le\ell\le4$, the eigenvalue is
$F=P_\ell(\grah(X))+2\grah(Y)$, where the integral polynomials satisfy
$P_\ell(0)=0$ and $P_\ell(z)\equiv z^\ell\pmod2$. Explicitly, the four
polynomials are $z$, $z(6-z)$, $z((3-z)^2+1)$ and $z(4-z)[4-z(4-z)]+2z$;
these identities follow from $\grah(X^2)=\grah(X)(4-\grah(X))$ and
$\grah(X^3)=\grah(X)(3-\grah(X))^2$.

Put $L=2^j$, the smallest power of two at least $\ell$, and take
$h\ge\max(3,j+2)$. If $a\ne0$ then
$t=v_2(\grah(X))=2\cdot2^{v_2(a)}/e$. When $t\le1/L$, the odd coefficient
of $z^\ell$ gives $v_2(P_\ell(\grah(X)))=\ell t\le1$, since lower
coefficients are even and the constant term vanishes. For odd $a$ this
applies regardless of $b$; the $2e^2$ choices contribute $4\ell e$ to
$S_h$. With $a$ even and $b$ odd, each stratum $v_2(a)=k$ for
$1\le k\le h-2-j$ contributes $2\ell e$: there are $e/2^k$ choices of $a$
and $e$ choices of $b$ in that stratum. In the remaining $2L$ choices of
$a$, either $a=0$ or $t\ge2/L$. For the latter, the leading monomial has
valuation $\ell t\ge2\ell/L>1+2/e$, and the terms of smaller positive
degree have valuation at least $1+t>1+2/e$; higher terms also exceed this
bound. The strict bound for $2\ell/L$ follows at once for
$(\ell,L)=(1,1),(2,2),(3,4),(4,4)$ and the stated range of $h$. The
valuation of $F$ for all these $a$ is thus $1+2/e$, that of $2\grah(Y)$.
In total,
\[
 S_h=\bigl[2\ell(h-j)+2L\bigr]e+4L .
\]

For $\lambda=3$ use the exact identity
$XYF=8XY-(1+X)(1+Y)(1+XY)$. For $h\ge2$, exactly one of $a,b,a+b$ is even;
label this residue $t$. The two odd arguments each give
$v_2(1+\zeta^{\mathrm{odd}})=1/e$. For each choice of the even position and
its value there are $e$ pairs. When $t=e$ the product is zero, hence
$v_2(F)=3$. For all the other even $t$ we have $v_2(1+\zeta^t)\le1$, hence
$v_2(F)=2/e+v_2(1+\zeta^t)<3$. The sum of $v_2(1+\zeta^t)$ over those
remaining even $t$ is $h-1$, since $-\zeta^t$ takes all nonidentity values
in $\gramu_e$ once and $\prod_{\xi\in\gramu_e,\,\xi\ne1}(1-\xi)=e$. Thus
$S_h=3(h+4)e-6$.

In the order $\lambda=2,4,6,8,3$, the shell totals at $h=1,2,3$ are
$(7,18,36)$, $(10,30,56)$, $(10,28,88)$, $(10,28,80)$ and $(9,30,78)$,
obtained from \cref{gra:lem:rules}(iii). Substituting in the recurrence and
then using the shell formulas gives the corresponding five rows.
\end{proof}

\begin{proof}[Proof for $\lambda\in\{5,7\}$]
The first list contains a basis, and the second contains $(2,1),(1,1)$,
another basis; thus all layers are connected. Put $\delta=\lambda-2$.
Direct multiplication gives
\[
 \overline F\doteq(X+1)(Y+1)(XY+1)^\delta\quad\text{in }
 \F_2[X^{\pm1},Y^{\pm1}],
\]
where $\doteq$ means equality up to a Laurent monomial. For each primitive
pair $(a,b)$ modulo $2^h$, exactly one of $a,b,a+b$ is even. Write $c$ for
this even member and $W=\zeta^c$, and let $2^t$ be the order of $W$,
including $t=0$ for $W=1$. For $t\ge1$ the order of $(1+z)^c-1$ in
characteristic two is $B=e/2^{t-1}$, taking $c$ to be its residue in
$[0,2^h)$; replacing an exponent by a congruent one does not change the
formal order below $2^h$. At an axis ($a$ even or $b$ even) the contact
order is $B+\delta+1$; at the diagonal ($a+b$ even) it is $2+\delta B$. If
the contact order is strictly below $e$, it equals $\ord_hF(\zeta^a,\zeta^b)$
by \cref{gra:lem:rules}(i). For $h\ge5$ this settles both axes when
$t\ge2$, and the diagonal when $t\ge L+1$, where $L=2$ for $\delta=3$ and
$L=3$ for $\delta=5$.

For the remaining classes let $\eta=(\lambda-1)/2$ and let $Z$ denote a
primitive root of order $2^h$. In an axis class expand
$W^\eta Z^\eta F(W,Z)$ as $\sum_jc_j(Z-1)^j$; symmetry treats the second
axis. In a diagonal class use $W^\eta ZF(Z,W/Z)$ instead. The following
table gives the smallest pair $(v_2(c_j),j)$ in lexicographic order; a dash
means that the generic calculation applies.
\[
\begin{array}{c|cc|cc}
&\multicolumn{2}{c|}{\text{axis}}&\multicolumn{2}{c}{\text{diagonal}}\\
 t&\lambda=5&\lambda=7&\lambda=5&\lambda=7\\\hline
 0&(1,2)&(1,4)&(2,2)&(2,2)\\
 1&(1,2)&(1,4)&(3,0)&(4,0)\\
 2&-&-&(2,0)&(3/2,2)\\
 3&-&-&-&(2,0)
\end{array}
\]
To verify the entries, at $W=1,-1$ respectively the axis coefficient
lists, in increasing $j$, are
\[
\begin{array}{c|ll}
5&(0,0,-10,-10,-2)&(8,16,6,-2,0)\\
7&(0,0,-28,-56,-42,-14,-2)&(-16,-48,-56,-32,-10,-2,0).
\end{array}
\]
For the diagonal the lists in increasing $t$ are
\[
\begin{array}{c|l}
5&(0,0,-4),\quad(8,8,0),\quad(-12,-12,0)\\
7&(0,0,-4),\quad(-16,-16,0),\quad
 (-16W,-16W-4,-2W-2),\quad(12W^3,12W^3,0).
\end{array}
\]
These expansions use only the identity $\Phi_{2^t}(W)=0$, and
$\Phi_1(W)=W-1$ when $t=0$. For $W$ of order four we have
$v_2(1+W)=1/2$, explaining the fractional entry. If the table entry is
$(H,J)$, the normalised order is $eH+J$: the coefficient valuations lie in
$\varphi(2^t)^{-1}\Z$, and for $h\ge5$ the polynomial degree (at most six
on the axes and two on the diagonal) is smaller than $e/\varphi(2^t)$ in
every tabulated case, so the minimum term is unique.

To sum each class over its primitive pairs, note that for each fixed $W$
there are $e$ choices for the primitive coordinate $Z$. A normalised order
$D$ therefore contributes $D$ for each such $W$, and there are
$\varphi(2^t)$ choices of exact order $2^t$. In a single axis class the
generic sum is $(h-2)e+(\delta+1)(e-2)$, while its two exceptional classes
contribute $2(e+\delta-1)$; each axis total is consequently
$(h+\delta+1)e-4$. The generic diagonal sum is
\[
 \delta(h-L-1)e+2(e-2^L),
\]
and the exceptional diagonal classes contribute $9e+2$ for $\delta=3$ and
$17e+6$ for $\delta=5$. Adding all three classes yields
\[
 S_h=(5h+10)e-14\quad(\lambda=5),\qquad
 S_h=(7h+11)e-18\quad(\lambda=7),\qquad h\ge5.
\]
The four initial values $(\nu_1,\dots,\nu_4)$ are $(8,34,134,358)$ and
$(10,50,154,446)$ respectively, by the finite cyclotomic basis calculation
of \cref{gra:lem:rules}(iii). Summation proves the asserted formulas.
\end{proof}

\begin{example}[A two-vertex base at every prime]\label[example]{gra:ex:honeycomb}
Let the base graph have two vertices joined by three edges with voltages
$0,e_1,e_2$; its covers are the periodic honeycomb graphs. The base has
three spanning trees, and the determinant of its two-by-two Laurent
Laplacian is
\[
 F(X,Y)=9-(1+X+Y)(1+X^{-1}+Y^{-1})=\grah(X)+\grah(Y)+\grah(X/Y).
\]
The file \artifact{Effective_Graph_Towers/honeycomb.json} encodes this
base. At $p=2$ and $p=3$ the algorithm of \cref{gra:thm:effective} yields
the all-level formulas
\[
 \nu_n=(3n+9)2^n-8n-9,\qquad\nu_n=5\cdot3^n-2n-4,
\]
with zero finite correction in both cases and register dimensions $44$
and $12$; the ternary constant term takes account of the base tree count.
For $p\ge5$, at every level,
\[
 \nu_n=\graA_p\,\frac{p^n-1}{p-1}-2n,\qquad
 \graA_p=\begin{cases}2(p+5),&p\equiv1\pmod3,\\2(p+1),&p\equiv2\pmod3.
 \end{cases}
\]
Indeed, in a primitive direction $(a,b)$ modulo $p$ the first even moment
is $a^2+b^2+(a-b)^2=2(a^2-ab+b^2)$. If it is nonzero, the first unit degree
is two. Otherwise $ab\ne0$ and $b/a$ is a primitive sixth root of unity;
the three squared projections $a^2,b^2,(a-b)^2$ are $a^2$ times the three
cube roots of unity, their first two power sums vanish and their third is
$3a^6\ne0$. There are two such projective directions for $p\equiv1\bmod3$
and none for $p\equiv2\bmod3$. To identify the first unit degree in the
exceptional directions, use $X=\exp(az)$, $Y=\exp(bz)$ formally through
degree six modulo $p$; here $p\ge7$, so all factorials in this expansion
are units. The coefficient at even degree $2j$ is $-2/(2j)!$ times the
$j$th power sum of the squared projections, and the change from $z$ to
$\exp(z)-1$ preserves the order, so the first unit degree is six. In a
cyclotomic specialisation the integral Laurent expansion in $T=\zeta-1$
has its first two coefficients identically zero; thus even at $p=7$, $h=1$,
when degree six equals the ramification index, every lower nonzero term
divisible by $p$ has normalised order at least $e+2>6$. These unit degrees
are exactly the orbit orders at every conductor. Each projective direction
has $p^{h-1}$ normalised lifts at level $h$; summing gives $\graA_p$ at
level one and the formula by geometric summation. This is the method of
\cref{gra:thm:moments} applied to a base with two vertices.
\end{example}

\subsection{Residue formulas when the prime exceeds the number of steps}\label{gra:sec:moments}
{\raggedright
\sources{grh_method_transfers_graph_towers_and_finite_rings; grh_transfers_lubin_tate_and_graph_towers; grh_transfer_extensions_2026-11-26; spanning-tree-sharp-cutoff}
}

The preceding arguments have a uniform form when the prime exceeds the
number of steps. Fix an odd prime $p$ and a bouquet with $m$ steps
$s_i\in\Z^d$ spanning $\F_p^d$ modulo $p$, with multiplicities kept. For
$a\in\F_p^d\setminus\{0\}$ define the following polynomial in
$\F_p[T]/(T^p)$:
\[
 L_a(T)=2m-\sum_{i=1}^m
 \bigl((1+T)^{[a\cdot s_i]}+(1+T)^{[-a\cdot s_i]}\bigr),
\]
where brackets denote representatives in $\{0,\dots,p-1\}$. When
$L_a\ne0$, let $j(a)$ be its least nonzero degree.

In the variable $Z=1+T$, which identifies the polynomials of degree less
than $p$ on both sides, the coefficient of $Z^r$ for $r\ne0$ is the
negative of the number of steps whose projection $a\cdot s_i$ lies in
$\{r,-r\}$, and the coefficients sum to zero. Hence $L_a\ne0$ precisely
when some nonzero unordered pair $\{r,-r\}$ of projected residues occurs
among the $m$ projected steps a number of times not divisible by $p$.
Spanning gives a nonzero projected residue in every nonzero direction, so
$m<p$ suffices for $L_a\ne0$ for all $a$.

\begin{theorem}[Residue formula below the sharp cutoff; \stU]\label{gra:thm:cutoff}
Let $p$ be odd, and suppose that $L_a\ne0$ for every
$a\in\F_p^d\setminus\{0\}$; this holds whenever $m<p$. Then for every
$n\ge1$
\[
 \nu_n=C\sum_{h=1}^np^{(d-1)(h-1)}-dn,\qquad
 C=\frac1{p-1}\sum_{a\in\F_p^d\setminus\{0\}}j(a).
\]
In particular, for fewer than $p$ steps the reductions of the steps modulo
$p$ determine the valuation at every level. The bound $m<p$ is sharp
(\cref{gra:prop:cutoffsharp}).
\end{theorem}

\begin{proof}
A character of exact order $p^h$ has a primitive vector
$b\in(\Z/p^h\Z)^d$. With $\zeta=\zeta_h$, $\pi=\zeta-1$ and $e=e_h$, its
Laplacian eigenvalue is $F_b(\zeta)$, where the integer polynomial $F_b(Z)$
is given by the formula for $L_a$ with $Z$ in place of $1+T$ and with
representatives modulo $p^h$. In particular $F_b(1)=0$. Over $\F_p$ one
has $(1+T)^p=1+T^p$, so the coefficients of $F_b(1+T)$ in degrees below
$p$, reduced modulo $p$, are those of $L_{\bar b}(T)$. Whenever
$j(\bar b)<e$, \cref{gra:lem:rules}(i) gives
$v_p(F_b(\zeta))=j(\bar b)/e$.

Since $j(\bar b)\le p-1\le e$, the only boundary case is $h=1$,
$j(\bar b)=p-1$. Then $F_b(Z)$ has degree at most $p-1$ and is congruent
to $c\,\Phi_p(Z)$ modulo $p$ for an integer $c$ prime to $p$, because
$\Phi_p(Z)\equiv(Z-1)^{p-1}\pmod p$. Write $F_b=c\,\Phi_p+pG$ with
$G\in\Z[Z]$. Evaluating at $1$ gives $G(1)=-c$. As $\zeta\equiv1$ modulo
$\pi$, $G(\zeta)$ is a unit, and $v_p(F_b(\zeta))=1=j(\bar b)/e_1$. This
uses the exact zero at $1$ and needs no bound on the multiplicities.

Each nonzero $a\in\F_p^d$ has $p^{d(h-1)}$ primitive lifts to
$(\Z/p^h\Z)^d$, so the shell sum is
$S_h=p^{d(h-1)}e_h^{-1}\sum_aj(a)=p^{(d-1)(h-1)}C$. The steps span, so
$\nu_0=0$, and \eqref{gra:eq:shells} gives the formula.
\end{proof}

\begin{remark}[More residue digits]\label[remark]{gra:rem:digits}
The same reduction works with more digits. If a primitive projection $b$
is known modulo $p^k$, expand $F_b(1+T)$ in $\F_p[T]/(T^{p^k})$ using
representatives modulo $p^k$. If its order there is $j<e_h$ and $h\ge k$,
all lifts to order $p^h$ give valuation $j/e_h$; in particular any nonzero
such polynomial suffices for $h>k$. This is useful when $L_{\bar b}=0$;
compare the reduction of the state space in \cref{gra:sec:effective}.
\end{remark}

When $p>2m$ the integers $j(a)$ are read off from power sums. For
$u\in\F_p^d$ and $j\ge1$ put
\[
 M_j(u)=\sum_{i=1}^m\langle u,s_i\rangle^{2j}\in\F_p ,
\]
a form of degree $2j$ in $u$. For a projective covector
$L\in\graP^{d-1}(\F_p)$ represented by $u$, define the \emph{depth} and the
\emph{moment coefficient}
\[
 r(L)=\min\{j\ge1:M_j(u)\ne0\},\qquad\graA=\graA(S)=2\sum_Lr(L).
\]
Scaling $u$ preserves the first nonvanishing degree.

\begin{theorem}[Moment formula; \stU]\label{gra:thm:moments}
Assume $p>2m$. Each $r(L)$ is at most $m$, and $j(a)=2r(L)$ for every
nonzero $a$ on $L$. Consequently the constant $C$ of \cref{gra:thm:cutoff}
equals $\graA$, and the formula there reads, for every $n\ge1$,
\[
 \nu_n=\graA\sum_{h=0}^{n-1}p^{(d-1)h}-dn .
\]
\end{theorem}

Thus under $p>2m$ the moment formula refines \cref{gra:thm:cutoff}, which
needs only $m<p$, by identifying its exponents $j(a)$ with twice the depths.

\begin{proof}
Put $x_i=\langle u,s_i\rangle^2$ in $\F_p$. If their first $m$ power sums
vanished, Newton's identities (with $1,\dots,m$ invertible) would give
$\prod_i(Z-x_i)=Z^m$. This forces all projections to be zero and
contradicts spanning. Next take the truncated exponential coordinate
$1+T=\exp z$ modulo $z^{2m+1}$ over $\F_p$; all denominators used are
units, and the linear term is invertible. Since the representative of
$-a\cdot s_i$ is congruent to $-[a\cdot s_i]$, the polynomial $L_a$ becomes
\[
 -2\sum_{j=1}^m\frac{z^{2j}}{(2j)!}\,M_j(a)\pmod{z^{2m+1}} .
\]
Its first nonzero degree is $2r(L)\le2m<p$, so $j(a)=2r(L)$. Each
projective point carries $p-1$ nonzero vectors, hence the constant $C$ of
\cref{gra:thm:cutoff} equals $\graA$, and that theorem gives the formula.
\end{proof}

\begin{example}\label[example]{gra:ex:torus}
For the square torus ($e_1,e_2$) and $p\ge5$, the first moment is
$a^2+b^2$. For $p\equiv3\pmod4$ it never vanishes on a projective point,
by the anisotropy behind \cref{gra:eq:aniso};
for $p\equiv1\pmod4$ its two isotropic directions have nonzero fourth
moment $2a^4$. Consequently $\graA=2(p+1)$ and $\graA=2(p+3)$,
respectively. At $p=5$ this gives $\nu_n=4(5^n-1)-2n$.
\end{example}

In rank one, Dion, Lei, Ray and Valli\`eres had already used the even
moments and analysed the distribution of the one-variable Iwasawa
invariant, including the two-loop case
\cite[Theorem~4.13 and Proposition~4.19]{DionLeiRayVallieres2024Distribution}.
\Cref{gra:thm:moments} counts all residual projective directions and their
lifts in arbitrary rank, and \cref{gra:thm:cutoff} removes the factor two
from its hypothesis $2m<p$.

\begin{proposition}[The cutoff $m<p$ is sharp; \stU]\label[proposition]{gra:prop:cutoffsharp}
Let $p$ be odd and $d=2$. For $0\le i<p$ take the $p$ steps $s_i=(1,i)$,
and alternatively the $p$ steps $s_i'=(1+pi^2,i)$. Both give simple
connected graphs at every level, and the two step lists agree modulo $p$.
Put $E=p(p-1)$ and $\delta=1$ if $p\equiv3\pmod4$, $\delta=0$ otherwise.
For every $n\ge1$ the two valuations are
\begin{align}
 \nu_n&=\sum_{j=0}^{n-1}\bigl[p^2-1+j(p-1)^2\bigr]p^j
 =\bigl((p-1)n+1\bigr)p^n-1,\label{gra:eq:flat}\\
 \nu_n'&=E+p+1+\Bigl(\frac{3E}2+\delta\Bigr)\sum_{j=1}^{n-1}p^j-2n .
 \label{gra:eq:curved}
\end{align}
They agree at $n=1$, and their difference at $n=2$ is
$pE/2+2-p\delta>0$. Thus for $m=p$ the residues of the steps determine
neither the valuations nor their growth order, which is $np^n$ for the
first list and $p^n$ for the second.
\end{proposition}

\begin{proof}
The steps and their negatives are distinct modulo $p$, and the first two
steps span; this gives simplicity and connectedness. Write the character
vector as $(a,b)$ at exact order $p^h$ and keep the notation of the
previous proof.

\emph{Common contributions.} At level one an order-$p$ character with
$b\ne0$ has eigenvalue $2p$, since the projections $a+bi$ run through all
residues. Otherwise the eigenvalue is $p(2-\zeta^a-\zeta^{-a})$, of
valuation $1+2/(p-1)$. The level-one contribution, before the vertex-count
term, is $E+p+1$ for both lists. For $h\ge2$ and $b$ a unit, the residue
polynomial is $-2\sum_{r=0}^{p-1}(1+T)^r=-2T^{p-1}$. Each such eigenvalue
has valuation $(p-1)/e_h$, and these characters contribute $Ep^{h-1}$ in
total. The remaining primitive vectors have $a$ a unit and $p\mid b$.

\emph{Straight steps.} If $b=0$ at order $p^h$, the contribution is
$e_h+2$. When $t=v_p(b)$ satisfies $1\le t\le h-1$, the eigenvalue is
\[
 2p-\bigl(\zeta^a+\zeta^{-a-(p-1)b}\bigr)\sum_{i=0}^{p-1}\zeta^{bi}.
\]
The parenthesis is a unit because it reduces to $2$ modulo $\pi$. If
$t<h-1$, the geometric sum has valuation $(p-1)p^t/e_h<1$, by writing it
as $(1-\zeta^{pb})/(1-\zeta^b)$. For $t=h-1$ the geometric sum is zero and
the eigenvalue has valuation $1$, which is again $(p-1)p^t/e_h$. There are
$e_h$ choices for $a$ and $(p-1)p^{h-1-t}$ choices for $b$ at each $t$, so
each $t$ contributes $(p-1)^2p^{h-1}$. Adding the common contributions and
subtracting $2n$ gives \eqref{gra:eq:flat}.

\emph{Quadratic steps.} Write $a=a_0+pa_1$ with $a_0\in\{1,\dots,p-1\}$,
and $b=pc$. The positive projected residues are congruent modulo $p^2$ to
$a_0+pk_i$, where $k_i=a_1+ci+a_0i^2\bmod p$. With $U=T^p$ and $Z=1+T$,
the sum of the positive projected terms modulo $(p,T^{p^2})$ is
$Z^{a_0}\sum_i(1+U)^{k_i}$. The negative sum has the factor $Z^{p-a_0}$ and
exponents $-1-k_i\bmod p$. Let $r=(p-1)/2$. For each $j<r$ the coefficient
of $U^j$ vanishes in both sums: it is a polynomial in $i$ of degree
$2j<p-1$, and every monomial of such degree, including the constant, sums
to zero over $\F_p$. At $j=r$ the positive coefficient is
$A=-a_0^r/r!\ne0$, coming from the term of degree $p-1$, and the negative
coefficient is $(-1)^rA$.

Consequently the first surviving degree in the Laplacian is $pr$ if $r$ is
even. If $r$ is odd, the constant of the $Z$-factor cancels, but its
coefficient of $T$ is $2a_0A$ up to the overall sign, which is nonzero, and
the degree is $pr+1$. In either case the degree is $E/2+\delta<e_h$, and
dividing by $e_h$ gives the eigenvalue valuation by
\cref{gra:lem:rules}(i). There are $e_hp^{h-1}$ such vectors, so the common
and quadratic contributions sum to $(3E/2+\delta)p^{h-1}$ for $h\ge2$.
Formula \eqref{gra:eq:shells} gives \eqref{gra:eq:curved}.
\end{proof}

At $p=3$, $n=2$ the two valuations are $44$ and $36$; independent integer
cofactors of the $80\times80$ reduced Laplacians give the same values
(\artifact{spanning-tree-sharp-cutoff/verification.txt}). In rank two the
pair of \cref{gra:prop:cutoffsharp} differs in the coefficient of $np^n$,
in accordance with the linear range $m\ge\max(3,p)$ of
\cref{gra:thm:loops}.

\subsection{The largest moment coefficient in rank two}\label{gra:sec:extremal}
{\raggedright
\sources{grh_transfers_lubin_tate_and_graph_towers; grh_transfer_extensions_2026-11-26; grh_method_transfers_graph_towers_and_finite_rings; short_tower_distributions_2026-11-27; five_step_graph_towers_2026-11-26; five_step_sharpness_2026-11-27; seven_step_graph_towers_2026-11-27; five_step_refinement_2026-11-27; sparse-padic-roots-and-graph-towers}
}

How divisible by $p$ can the spanning-tree count of a connected abelian
Cayley graph be, if the number of steps is fixed? By
\cref{gra:thm:moments} this is the problem of maximising $\graA(S)$. In
this subsection $d=2$, $p>2m$, and the $m$ steps span $\F_p^2$ after
reduction; repetitions and steps with zero residue are allowed. The forms
$M_j$ are binary forms of degree $2j$ over $\F_p$.

The value $\graA(S)$ depends only on the \emph{row plane}
$W\subset\F_p^m$ of the $2\times m$ matrix whose columns are the step
residues: a change of basis of $\F_p^2$ gives precisely the other matrices
with the same row plane, and a direction $L$ corresponds to the point of
$\graP(W)$ given by the projections $(\langle u,s_i\rangle)_i$. Permuting
the steps and changing their signs act on $W$ by signed permutations of
the coordinates. Distinct nonzero step residues up to sign give simple
graphs at every level.

\begin{theorem}[First nonzero moment form and the interpolation bound; \stU]\label{gra:thm:extremal}
Let $t$ be the first index $j$ for which $M_j$ is a nonzero polynomial, and
put $s=\floor{m/2}$. Then $t\le s$ and
\begin{equation}\label{gra:eq:tbound}
 \graA(S)\le2t(p+1)+4t(m-t)\le2s(p+1)+4s(m-s).
\end{equation}
A direction of depth $m$ exists only if $2m\mid p-1$. If no direction has
depth above $D$, then $\graA(S)\le2t(p+1+2D-2t)$; in particular, if
$2m\nmid p-1$,
\begin{equation}\label{gra:eq:tboundnc}
 \graA(S)\le2t(p+1)+4t(m-1-t)\le2s(p+1)+4s(m-1-s).
\end{equation}
For even $m=2s$, equality $\graA(S)=m(p+1+m)$ in the outer bound of
\eqref{gra:eq:tbound} is attainable exactly at the primes with
$2m\mid p-1$. The maximising configuration has $M_1=0$ identically
whenever $m\ge4$.
\end{theorem}

\begin{proof}
The degree of $M_t$ is $2t\le2m<p$. We first prove $t\le s$. Discard the
steps with zero residue and let $m'$ be the remaining count. If a line
through the origin contains more than half of these steps, take a nonzero
linear functional annihilating the line. Its list of projections has at
most $\floor{m'/2}$ nonzero entries and at least one, by spanning. Newton's
identities applied to the shorter list show that some $M_j$ with
$j\le\floor{m'/2}$ is nonzero at this covector.

Otherwise, arrange the $2m'$ signed occurrences $\pm s_i$ in contiguous
blocks according to the line through the origin that they occupy. Every
block has size at most $m'$. Pair occurrence $j$ with occurrence $j+m'$
for $1\le j\le m'$. Each pair has two distinct radial directions, so its
affine line avoids the origin. Multiplying the equations of these $m'$
affine lines produces a polynomial $H$ of degree $m'$ that vanishes on
every signed occurrence and satisfies $H(0)\ne0$. If $t>\floor{m'/2}$,
then $m'\le2t-1$. For every polynomial $P$ of degree at most $2t-1$ we
have
\[
 \sum_{i:s_i\ne0}\bigl(P(s_i)+P(-s_i)\bigr)=2m'P(0):
\]
odd homogeneous degrees cancel, and for degree $2j<2t$, $j\ge1$, the
vanishing coefficients of $M_j$ say that each mixed monomial of that degree
sums to zero, since all binomial coefficients $\binom{2j}{k}$ are units.
Applying the identity to $H$ contradicts $p>2m$. This proves $t\le s$ in
both cases.

A nonzero binary form of degree $2t$ has at most $2t$ projective roots.
The other directions have depth $t$, and each exceptional direction has
depth at most $D$, where $D\le m$ by \cref{gra:thm:moments}. Hence
$\graA/2\le t(p+1)+2t(D-t)$. For $D\in\{m-1,m\}$ the right side increases
strictly for integers $1\le t\le s$: the successive difference of
$2t(p+1+2D-2t)$ is $2(p+1+2D-4t-2)>0$ in that range because $p>2m$ and
$4t\le2m$.

In any direction of depth $m$, Newton's identities give
$\prod_i(X-\langle u,s_i\rangle^2)=X^m-c$ with $c\ne0$. The roots are
distinct nonzero squares. Dividing by one root shows that all $m$th roots
of unity belong to the square subgroup, and hence $2m\mid p-1$. This
proves \eqref{gra:eq:tbound} and \eqref{gra:eq:tboundnc}. For even $m=2s$,
equality in the outer bound of \eqref{gra:eq:tbound} requires $t=s$ and
$2s$ exceptional directions of depth $m$, hence $2m\mid p-1$.

Conversely, assume $4s\mid p-1$ and choose $g\in\F_p^\times$ of order
$2s$. Use the $s$ steps $(g^j,0)$ and the $s$ steps $(0,g^j)$,
$0\le j<s$, lifting residues to integers. Since $g^2$ has order $s$,
\[
 M_k(a,b)=\begin{cases}0,&s\nmid k,\\ s(a^{2k}+b^{2k}),&s\mid k.\end{cases}
\]
Exactly $2s$ rational directions solve $a^{2s}=-b^{2s}$, because $-1$ is a
$2s$th power; all have both coordinates nonzero. In each of them
$M_{2s}=2sb^{4s}\ne0$, so their depth is $2s=m$, and all other directions
have depth $s$. Thus
$\graA=2s(p+1-2s)+4s\cdot2s=m(p+1+m)$. For $s\ge2$ the display gives
$M_1=0$ identically.
\end{proof}

\begin{corollary}[Nonzero quadratic moment; \stU]\label[corollary]{gra:cor:quadratic}
If $M_1$ is not the zero polynomial, then $\graA(S)\le2(p-1)+4m$. If
$m\ge2$ and $2m\mid p-1$, this value is attained: take $g$ of order $2m$
modulo $p$ and lift the $m$ steps $(g^j,g^{-j})$, $0\le j<m$, to integers.
\end{corollary}

\begin{proof}
The bound is the case $t=1$ of \eqref{gra:eq:tbound}: a nonzero binary
quadratic form has at most two projective roots, each of which can
contribute at most $2m$. For the construction, the first two steps span.
In the direction $(a,b)$ the quadratic moment is $2mab$, since the sums of
$g^{2j}$ and of $g^{-2j}$ vanish. Each coordinate axis has all even
moments through order $2m-2$ zero, and its $2m$th moment is nonzero. Thus
$\graA=2(p-1)+4m$.
\end{proof}

For example $p=13$, $m=6$ gives $\graA=48$ in \cref{gra:cor:quadratic},
whereas the unrestricted maximum for the same prime and step count is
$120$:

\begin{corollary}[Exact maximum at every layer; \stU]\label[corollary]{gra:cor:layers}
For even $m$ and primes $p\equiv1\pmod{2m}$,
\[
 \max_Sv_p(\kappa_n(S))=m(p+1+m)\frac{p^n-1}{p-1}-2n\qquad(n\ge1),
\]
the maximum being over all lists of $m$ integral steps spanning $\F_p^2$.
The same fixed root-of-unity step construction attains the maximum
simultaneously at every level, independently of the chosen integral lifts.
\end{corollary}

\begin{proof}
Combine \cref{gra:thm:moments} with \cref{gra:thm:extremal}; the
congruence gives $p>2m$.
\end{proof}

\begin{example}\label[example]{gra:ex:thirteen}
For $m=6$, $p=13$ the maximum coefficient is $120$ and
$\max_Sv_{13}(\kappa_n(S))=10(13^n-1)-2n$. An attaining step list is
$(1,0),(4,0),(3,0),(0,1),(0,4),(0,3)$: the eight ordinary directions have
depth three and the six exceptional directions have depth six. Its values
at the first three levels are $118$, $1676$, $21954$. At $m=4$, $p=17$
the maximum coefficient is $88$ and the first-level valuation is $86$.
The value $118$ is recovered by evaluating the $168\times168$ reduced
integral Laplacian determinant directly; see
\grafiles{\artifact{grh_transfers_lubin_tate_and_graph_towers/checks/graph_checks.txt}\\
\artifact{grh_transfer_extensions_2026-11-26/verification/graph_check.txt}}
\end{example}

\paragraph{Three steps.}
For $m=3$ the bound \eqref{gra:eq:tbound} is $2(p+1)+8$; it is attained
when $p\equiv1\pmod6$ by \cref{gra:cor:quadratic}, and \eqref{gra:eq:tboundnc}
gives $2(p+1)+4$ otherwise. The full distribution is as follows.

\begin{theorem}[Three-step distribution; \stU]\label{gra:thm:three}
Let $m=3$ and $p>6$, put $E=\graA-2(p+1)$ and $t=p-7$. The numbers of row
planes $W\subset\F_p^3$ with each value of $E$ are:
\begin{center}\small\upshape
\begin{tabular}{crr}
\toprule
$E$ & $p\equiv1\pmod6$ & $p\equiv5\pmod6$\\
\midrule
0 & $p(p-1)/2$ & $p(p-1)/2$\\
2 & $t$ & $p+1$\\
4 & $8+\binom t2$ & $p(p+1)/2$\\
6 & $8t$ & $0$\\
8 & $28$ & $0$\\
\bottomrule
\end{tabular}
\end{center}
For the planes giving simple towers remove $3$ planes from $E=4$ if $-1$
is a square modulo $p$ and from $E=0$ otherwise, and remove $6$ planes
from $E=4$ if $-2$ is a square and from $E=0$ otherwise. In particular the
maximum of $\graA$ over three-step lists is $2(p+1)+8$ for
$p\equiv1\pmod6$ and $2(p+1)+4$ for $p\equiv5\pmod6$, also among simple
towers. For $p\equiv1\pmod6$ the $28$ maximising planes form exactly three
orbits under signed permutations, of sizes $12$, $4$ and $12$.
\end{theorem}

\begin{proof}
Every $\graP(W)$ is a line in $\graP^2$ and cannot be contained in the
smooth conic $Q:\sum x_i^2=0$. The excess $E$ is two for each point of
$Q$ on the line and an extra two for each point with also $\sum x_i^4=0$;
call the latter \emph{special}. A nondegenerate ternary form over a
finite field is isotropic (compare the two sets of square values in a
diagonal equation), and projection from a rational point of $Q$ gives
$p+1$ points. Thus there are $p+1$ tangents, $\binom{p+1}2$ secants and
$p(p-1)/2$ external lines. Newton's identities show that the special
points have squared coordinates the roots of $T^3-c$, $c\ne0$. They exist
exactly when $p\equiv1\pmod6$, in which case there are eight: normalise
the first coordinate to $1$, then order the other two squared coordinates
and choose their signs independently. Counting tangents and unordered
secant pairs by the number of special points proves the table.

Nonsimple row planes have a normal proportional to a coordinate vector or
to $e_i\pm e_j$; there are three of the first type and six of the second.
Their lines impose $x_i=0$ or $x_i=\pm x_j$, so contain no special point.
Substituting into the conic gives a secant or an external line according
as $-1$ or $-2$ is a square or a nonsquare. Hence the simple-tower maximum
agrees with the unrestricted one in both congruence classes.

For $p\equiv1\pmod6$ fix $z$ of order three and the special point
$a=(1,z,z^2)$. The signed permutation group is transitive on the eight
special points. The stabiliser of $a$ contains the cyclic permutation of
the coordinates, acting projectively. On the other seven special points
it fixes $(1,z^2,z)$ and cycles the three points of each of the two
remaining types: a single sign change of $a$, or a single sign change of
$(1,z^2,z)$. Cross products give the normals
\[
 (0,-z,1),\qquad(1,1,1),\qquad(z^2-z,1,-1)
\]
of the maximising lines through $a$ of the respective types, the middle
one being the line through the fixed point. These types are inequivalent
even as unordered pairs: the first normal has a zero coordinate, the
second has three equal squares, and the third has squares $-3,1,1$. Hence
the $28$ maximising planes form exactly three orbits, of sizes $12,4,12$.
\end{proof}

Distributions at thirteen primes from $7$ to $103$, obtained by direct
incidence enumeration, agree with the table; see
\grafiles{\artifact{short_tower_distributions_2026-11-27/verification.txt}}

\begin{theorem}[Four steps at every prime; \stU]\label{gra:thm:four}
Let $m=4$ and let $p>8$ be prime. Let $C\subset\graP^3$ be the curve
\[
 C:\quad x_1^2+x_2^2+x_3^2+x_4^2=x_1^4+x_2^4+x_3^4+x_4^4=0,
\]
and put $c_p=\#C(\F_p)$.
\begin{enumerate}[label=\textup{(\alph*)}]
\item If $M_1$ is not the zero polynomial, then $\graA(S)\le2(p+1)+12<4(p+1)$.
\item\label{gra:thm:fourdist} Among the row planes with $M_1=0$, put
$E=\graA-4(p+1)$ and $k=12$ if $p\equiv1\pmod8$, $k=0$ otherwise. The
numbers of such row planes with $E=0,8,16$ are respectively
\[
 2(p+1)-k-c_p/2,\qquad c_p/2,\qquad k .
\]
All these planes give simple towers. If $p\equiv1\pmod8$, the twelve
maximising planes form one orbit under signed permutations; thus the
maximising residue lists form one orbit under change of basis, signs and
permutations.
\item If $p\not\equiv1\pmod8$, then $c_p=0$ if and only if
$p\in\{11,23,29,47,53,83,149\}$.
\end{enumerate}
Consequently the unrestricted maximum of $\graA(S)$, which is also the
maximum over lists with $M_1=0$, is
\begin{equation}\label{gra:eq:four}
 \graA_{\max}=4(p+1)+\begin{cases}
 16,&p\equiv1\pmod8,\\
 0,&p\in\{11,23,29,47,53,83,149\},\\
 8,&\text{otherwise}.\end{cases}
\end{equation}
\end{theorem}

\begin{proof}
\emph{Part \textup{(a)}.} If $M_1$ is nonzero, the at most two exceptional
directions give $\graA\le2(p+1)+12$, the case $t=1$ of
\eqref{gra:eq:tbound}, and $2(p+1)+12<4(p+1)$ since $p>8$.

\emph{The Klein four group.} Now let $M_1=0$. Relabel two independent steps
and make them $e_1,e_2$ by a change of basis. The two rows of the remaining
$2\times2$ block must be perpendicular, both with sum of squared
coordinates $-1$. The first row $(a,b)$ determines the second up to sign.
Changing the sign of the second basis coordinate, followed by a sign change
of the second step, if needed, puts the block in the form
\[
 B=\begin{pmatrix}a&b\\b&-a\end{pmatrix},\quad a^2+b^2=-1,\qquad
 J=\begin{pmatrix}0&-1\\1&0\end{pmatrix}.
\]
Such $a,b$ exist for every odd prime because the squares and $-1$ minus
the squares each have $(p+1)/2$ elements. \Cref{gra:thm:extremal} gives
$t=2$.

Both $B$ and $J$ permute the signed steps
$\{\pm e_i,\pm Be_i:i=1,2\}$, with multiplicities. They satisfy
$J^2=B^2=(JB)^2=-I$ and $JB=-BJ$. Their projective transformations on
covectors generate a group of order four with three involutions: $B$ is
symmetric and nonscalar, $J$ is skew-symmetric and invertible, and $JB$ is
projectively different from both. Every $M_j$ is invariant under the
corresponding transpose matrices. A direction fixed by any nonidentity
element has eigenvalue $\lambda$ with $\lambda^2=-1$ for that
representative, so $M_3(u)=\lambda^6M_3(u)=-M_3(u)$. If such a direction
were a root of the quartic $M_2$, it would have depth four, which requires
$8\mid p-1$ by \cref{gra:thm:extremal}. Thus on a row plane with no
direction of depth four, in particular on every row plane with $M_1=0$ when
$p\not\equiv1\pmod8$, all rational roots of $M_2$ are in orbits of four.
Since $M_2$ is a nonzero quartic, it has either zero or four rational
roots, and each has depth three. This contributes precisely zero or eight
above $4(p+1)$. The projections $x$ in a root direction satisfy
$\sum x_i^2=\sum x_i^4=0$, so the roots of $M_2$ on a row plane are its
points on $C$.

\emph{Part \textup{(b)}.} The quadric $\sum_{i=1}^4x_i^2=0$ is split:
choose $u^2+v^2=-1$ as above and take the orthogonal isotropic rows
$(1,0,u,v)$, $(0,1,-v,u)$. Its two rulings therefore give $2(p+1)$
rational planes, and exactly two of them pass through each rational point
of the quadric; these are the row planes with $M_1=0$. The quartic
restricts nontrivially to every such plane by \cref{gra:thm:extremal}.
Depth four requires that the squared coordinates form a coset of the fourth
roots of unity inside the nonzero squares. Let $D$ be the set of depth-four
points; it is empty unless $p\equiv1\pmod8$. Otherwise, with $g$ of order
eight, the points of $D$ are the signed permutations of $a=(1,g,g^2,g^3)$
up to scaling; normalising the first coordinate to $1$ gives
$3!\,2^3=48$ points.

Exactly two isotropic planes pass through $[a]$, obtained by adjoining
$b=(-g,1,-g^3,g^2)$ and $d=(-g,1,g^3,-g^2)$ respectively; they are
distinct since $b\cdot d=2(1+g^2)\ne0$. On the first plane the coordinate
linear forms are $U,V,g^2U,g^2V$ with $U=X-gY$, $V=gX+Y$. The quartic is
$2(U^4+V^4)$, with four distinct rational roots, and the sextic vanishes
identically, so all four roots lie in $D$. On the second plane the quartic
is $24g^2X^2Y^2$; both axis points lie in $D$, and there are no other
roots. Consequently a plane meeting $D$ has either four points of $D$ and
excess $16$, or two points of $D$ and excess $8$. Counting incidences with
the $48$ points of $D$ gives twelve planes of the first type and
twenty-four of the second. Transitivity on $D$ and uniqueness of the
four-point plane through each point of $D$ give the orbit assertion.

On a plane avoiding $D$ the quartic has zero or four rational roots, each
of depth three, by the Klein four group argument above; such a plane has
excess $0$ or $8$. Every point of $C\setminus D$ lies on two planes
avoiding $D$. Hence $(c_p-\#D)/2$ planes avoiding $D$ have excess $8$, in
addition to the $\#D/2$ planes with two points of $D$, which proves the
counts. Simplicity is \cref{gra:rem:simple} below.

\emph{Part \textup{(c)}.} Over an algebraic closure $C$ is a smooth
complete intersection: dependence of its two gradients would force all
nonzero $x_i^2$ to coincide with one common nonzero $c$, and the quadratic
equation would give $kc=0$ for some $1\le k\le4<p$. The quartic does not
vanish identically on the smooth quadric, as the point $(1,\sqrt{-1},0,0)$
shows. On this quadric, isomorphic over the algebraic closure to
$\graP^1\times\graP^1$, the curve has bidegree $(4,4)$. It is connected:
if two nonempty effective divisors of classes $(a,b)$ and $(4-a,4-b)$
partitioned it disjointly, their intersection number $a(4-b)+b(4-a)$ would
be positive. Smoothness then gives geometric irreducibility, and
adjunction gives genus $(4-1)(4-1)=9$. The Hasse--Weil theorem
\cite{Weil1948Courbes} yields
\[
 \#C(\F_p)\ge p+1-18\sqrt p>0\qquad(p>322).
\]
For the $49$ remaining primes $8<p\le322$, $p\not\equiv1\pmod8$, normalise
a nonzero coordinate to one using permutation and scaling, and loop
through $y,z\in\F_p$: set $c=-1-y^2-z^2$ and test whether $c$ is a square
and $1+y^4+z^4+c^2=0$. This exhausts $C(\F_p)$ for the existence test.
A witness for every successful prime, and an independent count of the
whole curve at the $62$ primes from $11$ through $317$ by convolving
two-coordinate moment tables, are recorded in
\grafiles{\artifact{grh_transfers_lubin_tate_and_graph_towers/checks/graph_checks.txt}\\
\artifact{grh_transfers_lubin_tate_and_graph_towers/checks/graph_independent_output.txt}}
The seven failures are exactly the primes in \eqref{gra:eq:four}.

\emph{The maximum.} By (a), lists with $M_1\ne0$ are not maximisers. By
(b) the largest excess over $4(p+1)$ among lists with $M_1=0$ is $16$ if
$p\equiv1\pmod8$, $8$ if $p\not\equiv1\pmod8$ and $c_p>0$, and $0$ if
$c_p=0$. Part (c) turns this into \eqref{gra:eq:four}.
\end{proof}

\begin{remark}[Simple graphs]\label[remark]{gra:rem:simple}
The four-step maximum is also attained among simple eight-regular Cayley
graphs. The row plane $W\subset\F_p^4$ of a spanning list with $M_1=0$ is
totally isotropic for the dot product. A zero step, or equality of two
steps up to sign, would put $W$ in the orthogonal complement of $e_i$ or
of $e_i\pm e_j$. Each complement is a nondegenerate three-space, in which
a totally isotropic plane would violate $\dim W+\dim W^\perp=3$. Distinct
nonzero steps up to sign modulo $p$ give a simple graph at every level.
\end{remark}

\begin{remark}[Maximising row planes]
For $p\not\equiv1\pmod8$ the number of maximising row planes is $c_p/2$ if
$c_p>0$ and $2(p+1)$ if $c_p=0$; for $p\equiv1\pmod8$ it is twelve. The
extremal lists with $M_1=0$ are the maximal isotropic planes
$W\subset\F_p^4$ for the dot product. Reduced Laplacians at $p=11$ and
$p=13$ give four-step valuations $46$ and $62$ at level one, in agreement
with \eqref{gra:eq:four} and \cref{gra:thm:moments}. The distributions at
sixteen primes from $11$ to $509$, computed from both rulings with a
separate chart count of $C$, agree with
\cref{gra:thm:four}\ref{gra:thm:fourdist}:
\grafiles{\artifact{short_tower_distributions_2026-11-27/verification.txt}}
\end{remark}

\paragraph{The curve $C$ and the complete four-step census.}
The count $c_p$ is governed by two fixed elliptic curves over $\Q$, and
this gives the distribution of $\graA$ over all row planes, including those
with $M_1\ne0$. Put
\[
 E_1:\ y^2=x(x+2)(x-1),\qquad E_2:\ y^2=x(x-1)(x-4),
\]
and for $p\ge5$ let $a_p=p+1-\#E_1(\F_p)$ and $b_p=p+1-\#E_2(\F_p)$.

\begin{theorem}[The Jacobian of the moment curve; \stU]\label{gra:thm:jacobian}
Let $C\subset\graP^3$ be the curve $\sum_ix_i^2=\sum_ix_i^4=0$ of
\cref{gra:thm:four}. Then $\graJac(C)\sim_\Q E_1^6\times E_2^3$. For every
$p\ge5$ the zeta function of $C/\F_p$ has numerator
$(1-a_pT+pT^2)^6(1-b_pT+pT^2)^3$; in particular
\[
 c_p=\#C(\F_p)=p+1-6a_p-3b_p .
\]
\end{theorem}

\begin{proof}
The curve is smooth of genus nine in characteristic zero and at every
$p\ge5$, by the argument in the proof of \cref{gra:thm:four}(c). Put
$y_i=x_i^2$. The image is the conic $Y_0:\sum y_i=\sum y_i^2=0$, and the
cover $C\to Y_0$ has group $G=\{\pm1\}^4/\{\pm(1,1,1,1)\}$ of order eight,
as is seen on the open set with no coordinate zero. Its seven nontrivial
characters are indexed by the six pairs and the full set of four indices.
For an even subset $I$ the corresponding quotient double cover has equation
$t^2=\prod_{i\in I}y_i$ over $Y_0$, read homogeneously with $t$ of weight
$\abs I/2$. Each coordinate line meets $Y_0$ in two simple points, and the
four pairs are disjoint: with one coordinate zero, the other three are
proportional to the three cube roots of unity in some order, and two zero
coordinates are impossible. Riemann--Hurwitz gives genus one for each pair
cover and genus three for the full cover $H$.

Quotient differentials pull back to the subspace invariant under the
quotient subgroup, by descent for regular differentials on smooth curves
in characteristic zero. Since the $G$-invariant part is zero, the pair and
full-cover subspaces occupy distinct character spaces, and their
dimensions add to nine. The sum of the seven pullbacks on Jacobians is
therefore an isogeny defined over $\Q$, and rational permutations identify
the six pair covers.

In the chart $y_1+y_2=1$ let $v=y_1-y_2$ and $w=y_3-y_4$; then
$y_3+y_4=-1$ and $v^2+w^2=-2$. On the pair cover for $\{1,2\}$ take $h=2t$,
so $h^2=1-v^2$. The map $(v,w,h)\mapsto(v^2,vwh)$ takes values on $E_1$ and
is nonconstant; it extends to a morphism of smooth projective curves, and
the induced map on Jacobians is a rational isogeny, even if the source has
no rational base point. On $H$ take $h=4t$, so $h^2=(1-v^2)(v^2+3)$. The
double transposition $\sigma=(12)(34)$ negates $v,w$ and fixes $h$. With
$s=v^2+1$,
\[
 (vwh)^2=(s^2-1)(s^2-4),\qquad(v,w,h)\longmapsto(s^2,svwh)\in E_2 ,
\]
a nonconstant map factoring through $H/\langle\sigma\rangle$. On the conic,
in the coordinates $(y_1+y_2,y_1-y_2,y_3-y_4)$, $\sigma$ fixes exactly the
two points with $y_1+y_2=0$; both lie off the branch divisor, and the lift
fixes both points above each because $t$ has even weight. So $\sigma$ has
four fixed points on $H$ and Riemann--Hurwitz gives genus one for the
quotient. The three nonidentity elements of the double-transposition Klein
group $V$ are conjugate by rational permutations, so all three quotients
have genus one. If $m_0$ is the dimension of the $V$-invariant
differentials and $m_1,m_2,m_3$ those of the other characters, then
$m_0+m_i=1$ and $m_0+m_1+m_2+m_3=3$, so $m_0=0$ and $m_i=1$. The three
pullbacks give $\graJac(H)\sim_\Q E_2^3$. Good reduction of all three curves
at $p\ge5$ and invariance of the Frobenius polynomial under isogeny give
the remaining assertions.
\end{proof}

The Hasse bounds for $E_1$ and $E_2$ give $\abs{c_p-p-1}\le18\sqrt p$, the
bound used in the proof of \cref{gra:thm:four}(c). Substituting
$(y_i)=(l+m+n,l-m-n,-l+m-n,-l-m+n)$ turns the conic into
$l^2+m^2+n^2=0$, on which $\prod_iy_i=2(l^4+m^4+n^4)$; this identifies
$H$ with the central quadratic twist by $-1$ of the model of
\cite[Eq.~(1.4.1)]{Arora2016Twisting}, after rescaling the coordinate of
weight two, and the explicit map fixes the twist of $E_2$ in the trace
formula. The geometric multiplicity pattern $3+6$ for a genus-nine curve
with automorphism group of order $192$, and the decomposition of the
octahedral genus-three curve into an elliptic cube, are due to Paulhus
\cite[\S\S3.2.4, 4.2]{Paulhus2008Decomposing}; Arora et al.\ analysed the
rational twists of that genus-three curve \cite{Arora2016Twisting}. The
rational quotient maps above are what fix the signs of the twists.

\begin{corollary}[Diagonal twists; \stU]\label[corollary]{gra:cor:twists}
For nonzero $d_1,\dots,d_4\in\Q$ let $C_{\boldsymbol d}$ be the curve
$\sum_id_ix_i^2=\sum_id_i^2x_i^4=0$, and write $E^{(u)}$ for the quadratic
twist by $u$. Then
\[
 \graJac(C_{\boldsymbol d})\sim_\Q\prod_{i<j}E_1^{(d_id_j)}\times
 \bigl(E_2^{(d_1d_2d_3d_4)}\bigr)^3 .
\]
\end{corollary}

\begin{proof}
Over the field generated by the $\sqrt{d_i}$ the substitution
$x_i\mapsto\sqrt{d_i}x_i$ identifies the curves, and its descent action is
a coordinate sign change. On the factor for $I$ this acts by the product of
the signs indexed by $I$, since the deck involution acts by $-1$ on its
differentials, hence on its Jacobian (the quotient being the conic). This
is exactly twisting by $\prod_{i\in I}d_i$, and the same sign acts on all
three elliptic factors of $\graJac(H)$.
\end{proof}

For $p\ge5$ with all $d_i$ units, let $k$ be the number of nonsquare $d_i$
modulo $p$. The Frobenius trace on the twisted Jacobian is $6a_p+3b_p$,
$-3b_p$ or $-2a_p+3b_p$ according as $\min(k,4-k)=0,1,2$; in the middle
case the sixfold contribution disappears from the trace.

Put $\ell=p+1$, $\iota=\mathbf 1_{p\equiv1\pmod8}$, $c=c_p$ and
$\delta=48\iota$. The rational points of the split quadric
$Q:\sum_ix_i^2=0$ form an $\ell\times\ell$ grid, with $2\ell$ generator
lines, two through each point. Give the points of $Q\setminus C$,
$C\setminus D$ and $D$ (the depth-four points of the proof of
\cref{gra:thm:four}\ref{gra:thm:fourdist}) the weights $r-1=1,2,3$, and encode them by
\[
 F(z)=(\ell^2-c)z+(c-\delta)z^2+\delta z^3,\qquad
 \operatorname{Pairs}(h)(z)=\frac{h(z)^2-h(z^2)}2 ,
\]
so that $\operatorname{Pairs}$ counts unordered pairs of distinct weighted
points. For the four types of generator line let $f_j$ record its points,
$n_j$ its multiplicity and $e_j$ its total weight:
\begin{center}\small\upshape
\begin{tabular}{llrr}
\toprule
type $j$ & $f_j(z)$ & $n_j$ & $e_j$\\
\midrule
$0$ & $\ell z$ & $2\ell-12\iota-c/2$ & $\ell$\\
$B$ & $(\ell-4)z+4z^2$ & $(c-\delta)/2$ & $\ell+4$\\
$D_2$ & $(\ell-2)z+2z^3$ & $24\iota$ & $\ell+4$\\
$D_4$ & $(\ell-4)z+4z^3$ & $12\iota$ & $\ell+8$\\
\bottomrule
\end{tabular}
\end{center}

\begin{theorem}[Distribution of $\graA$ over all four-step row planes; \stU]\label{gra:thm:census}
For $p>8$ let $\Gamma_p(z)=\sum_Wz^{\graA(W)/2-\ell}$, the sum over all
two-dimensional $W\subset\F_p^4$. Then
\[
 \Gamma_p(z)=\frac{p^2(p-1)^2}2+(p-1)F(z)+\operatorname{Pairs}(F)(z)
 +\sum_jn_j\bigl(z^{e_j}-\operatorname{Pairs}(f_j)(z)\bigr).
\]
Thus the coefficient of $z^k$ is the number of row planes with
$\graA=2(p+1+k)$, and by \cref{gra:thm:jacobian} the only arithmetic input
is the pair $(a_p,b_p)$. Since every row plane has exactly
$\abs{\mathrm{GL}_2(\F_p)}$ ordered bases, this is also the law of $\graA$,
and by \cref{gra:thm:moments} of $v_p(\kappa_n)$ at every level, for a
uniformly random spanning residue matrix with four columns.
\end{theorem}

\begin{proof}
The generator table is the content of \cref{gra:thm:four}\ref{gra:thm:fourdist} and its
proof: the generators through $D$ are of types $D_4$ and $D_2$, each point
of $D$ lying on one of each, and the generators avoiding $D$ carry zero or
four points of $C$, each point of $C\setminus D$ lying on two of them. Each
row plane corresponds to a projective line, and its exponent in $\Gamma_p$
is the sum of $r(P)-1$ over its rational points; points off $Q$ have depth
one. On a line that is not a generator the sum comes only from its
intersection with $Q$. Every pair of distinct points of $Q$ spans one
secant unless its line is a generator, and no pair lies on two generators;
this gives $\operatorname{Pairs}(F)-\sum_jn_j\operatorname{Pairs}(f_j)$. The
tangent plane at a point of $Q$ contains $p+1$ lines through it, and
removing its two generators leaves $p-1$ tangents, which gives $(p-1)F$.
Each point of $Q$ is joined to $p^2$ others outside its two generators, so
there are $\ell^2p^2/2$ secants. Subtracting these, the $\ell^2(p-1)$
tangents and the $2\ell$ generators from the $(p^2+1)(p^2+p+1)$ lines of
$\graP^3$ leaves $p^2(p-1)^2/2$ external lines, of exponent zero. Adding
the generators themselves proves the formula.
\end{proof}

At $p=17$ one has $a_p=-6$, $b_p=2$, so $c=\delta=48$: every one of the
$36$ generators meets $D$, $24$ have $\graA=80$ and $12$ have $\graA=88$,
and the latter give $v_{17}(\kappa_n)=88(17^n-1)/16-2n$. At $p=13$,
$c=32$ and $\delta=0$, so $16$ of the $28$ generators have $\graA=64$ and
$12$ have $\graA=56$. As $p\to\infty$ the external lines account for a
proportion of all row planes tending to $1/2$, all with $\graA=2(p+1)$, and
the secants through two points of $Q\setminus C$ for another proportion
tending to $1/2$, all with $\graA=2(p+1)+4$; the remaining secants involve
one of the $O(p)$ points of $C$ and number $O(p^3)$.

\begin{theorem}[A congruence modulo $192$; \stU]\label{gra:thm:c192}
Let $J_p=1$ if $-2$ is a square modulo $p$ and both $-1\pm\sqrt{-2}$ are
squares in $\F_p$, and $J_p=0$ otherwise. For $p\ge5$,
\[
 c_p\equiv32\cdot\mathbf 1_{p\equiv1\pmod3}+48\cdot\mathbf 1_{p\equiv1\pmod8}
 +96J_p\pmod{192}.
\]
For every modulus $M$, even among the primes $p\equiv1\pmod{24M}$, the two
possible values $80$ and $176$ of $c_p$ modulo $192$ each occur with
relative density $1/2$.
\end{theorem}

\begin{proof}
The projective signed permutations, a group of order $8\cdot24=192$, act on
$C(\F_p)$. For $y_i=x_i^2$ the polynomial with these four roots is
$T^4-e_3T+e_4$, by the two equations and Newton's identities. If
$e_3e_4\ne0$, a rescaling preserving the multiset of the $y_i$ has cube and
fourth power equal to one, so it is the identity on the squares; when the
four squares are distinct and nonzero the projective stabiliser is trivial
and the orbit has size $192$. There are three disjoint exceptions. If
$e_4=0$, the squares are projectively $0,1,\omega,\omega^2$ with $\omega$ a
primitive cube root of unity: four choices of the zero, two orderings after
scaling a chosen nonzero coordinate to one, and four sign choices give $32$
points when $3\mid p-1$. If $e_3=0$, these are the $48\iota$ points of
depth four (the same roots-of-unity count works at $p=5,7$). In the
remaining case of a repeated root the squares scale to
$1,1,-1+\sqrt{-2},-1-\sqrt{-2}$, since the sum and the sum of squares force
the last two to be the roots of $T^2+2T+3$; these are distinct and differ
from $0,1$ for $p\ge5$. Six choices of the equal pair, two orders of the
remaining values and eight sign choices give $96J_p$.

For the density let $\alpha,\beta=-1\pm\sqrt{-2}$ and
$K=\Q(\sqrt{-2},\sqrt\alpha,\sqrt\beta)$, the splitting field of
$T^4+2T^2+3$. The classes of $\alpha$, $\beta$ are independent modulo
squares in $\Q(\sqrt{-2})$: their norms are $3$, and $\alpha\beta=3$ is not
a square there. Hence $\operatorname{Gal}(K/\Q)$ is dihedral of order
eight; conjugation exchanges the two square roots and the two independent
sign changes, and the commutator changes both signs, so the largest
abelian subfield is $L=\Q(\sqrt{-2},\sqrt3)\subset\Q(\zeta_{24})$. For
$p\ge5$, $J_p=1$ exactly when $p$ splits completely in $K$. With
$F=\Q(\zeta_{24M})$ one has $K\cap F=L$; in $KF$ the subgroup fixing $F$ has
order two and is central, being identified by restriction with the
commutator subgroup of $\operatorname{Gal}(K/\Q)$. The Chebotarev density
theorem gives equal densities to its two elements among the primes that
split in $F$, and the congruence gives the two residues.
\end{proof}

\begin{corollary}[Fluctuations of the middle class; \stU]\label[corollary]{gra:cor:fluct}
Let $N_p=c_p/2$ be the number of generators with $\graA=4(p+1)+8$, and
$T_p=(N_p-(p+1)/2)/\sqrt p$. As $p$ ranges over the primes $p>8$, $T_p$ has
limiting distribution $-3X-\frac32Y$, where $X,Y$ are independent with
density $\sqrt{4-x^2}/(2\pi)$ on $[-2,2]$; its limiting mean is $0$ and its
limiting variance $45/4$, and
\[
 \liminf_{p\to\infty}T_p=-9,\qquad\limsup_{p\to\infty}T_p=9 .
\]
The proportion of the $2(p+1)$ generators in this class tends to $1/4$, and
it is above $1/4$ on a set of primes of density $1/2$ and below it on a set
of density $1/2$.
\end{corollary}

\begin{proof}
By \cref{gra:thm:four}\ref{gra:thm:fourdist} and \cref{gra:thm:jacobian},
$T_p=-(6a_p+3b_p)/(2\sqrt p)$. The $j$-invariants of $E_1,E_2$ are
$21952/9$ and $35152/9$; they are not integers, so neither curve has
complex multiplication. At $p=7$ the traces are $-4$ and $0$, so $E_1$ is
ordinary and $E_2$ supersingular there, which excludes even a geometric
isogeny, since isogenies specialise at primes of good reduction and
preserve this distinction. Thus the two non-CM weight-two newforms are not
twist-equivalent, and the joint Sato--Tate theorem
\cite[Theorem~1.2]{Thorner2021Effective} gives independent semicircle
limits for $a_p/\sqrt p$ and $b_p/\sqrt p$. Each has mean zero and variance
one. The product measure gives positive mass arbitrarily near each corner
of its square, and the Hasse bounds give $\abs{T_p}\le9$. Symmetry and
continuity of the limit give the density statements.
\end{proof}

\begin{proposition}[Mean and variance in every rank; \stU]\label[proposition]{gra:prop:meanvar}
Let $p>8$, $1\le k\le4$, and let $W$ be a uniformly random
$k$-dimensional row space of four step columns, with
$\graA=2\sum_{P\in\graP(W)}r(P)$. Put $u_j=(p^j-1)/(p-1)$,
\[
 \rho=\frac{u_k}{u_4},\qquad\tau=\frac{u_ku_{k-1}}{u_4u_3},\qquad
 U=\ell^2+c+\delta,\qquad V=\ell^2+3c+5\delta .
\]
Then
\[
 \mathbb E[\graA]=2u_k+2\rho U,\qquad
 \operatorname{Var}(\graA)=4\bigl[(\rho-\tau)V+(\tau-\rho^2)U^2\bigr].
\]
\end{proposition}

\begin{proof}
By double counting and transitivity, a fixed point of $\graP^3$ lies in
$\graP(W)$ with probability $\rho$, and a fixed ordered pair of distinct
points with probability $\tau$. The sums over $\graP^3$ of the weights
$r(P)-1$ and of their squares are $U$ and $V$, by the three types of points
on $Q$. Expanding the first and second moments of
$\graA/2=u_k+\sum_{P\in\graP(W)}(r(P)-1)$ gives the formulas.
\end{proof}

\noindent The census is compared with direct row-space enumeration, and the
two elliptic character sums with point counts from the defining equations,
in
\grafiles{\artifact{sparse-padic-roots-and-graph-towers/graph/verification.txt}}

\paragraph{Five steps.}
For $m=5$ the bound \eqref{gra:eq:tbound} is $4(p+1)+24$.

\begin{theorem}[Five steps; \stU]\label{gra:thm:fivestep}
Let $m=5$ and $p>10$.
\begin{enumerate}[label=\textup{(\alph*)}]
\item $\graA(S)=4(p+1)+24$ for some list if and only if $p=11$. At $p=11$
every maximiser is equivalent to
\[
 \bigl((1,0),(0,1),(0,2),(1,1),(3,7)\bigr)
\]
under a common invertible linear map on the residues, permutation and
independent signs. This list gives simple ten-regular graphs, and for
every $n\ge1$
\[
 \max_Sv_{11}(\kappa_n(S))=\frac{72(11^n-1)}{10}-2n .
\]
\item At $p=31$ the maximum is $148=4(p+1)+20$, attained by one orbit,
represented by $((1,1),(2,29),(4,8),(8,16),(16,27))$.
\item For every other prime $p>10$, $\graA(S)\le4(p+1)+16$.
\item Let $p\ne11,31$. If exactly one direction has depth five, then
$\graA(S)\le4(p+1)+14$. This value is attained for infinitely many primes
$p\equiv1\pmod{10}$ by lists with exactly one direction of depth five whose
steps are nonzero and distinct up to sign, hence by simple towers.
\item If $p\equiv1\pmod{10}$ and $\xi$ has order five modulo $p$, the steps
$(\xi^i,\xi^{2i})$, $0\le i<5$, give $\graA=4(p+1)+12$.
\end{enumerate}
\end{theorem}

\begin{proof}
By \cref{gra:thm:moments} every depth is at most five. If $M_1\ne0$, then
$\graA\le2(p+1)+16$ by \eqref{gra:eq:tbound} with $t=1$, below all
values claimed. Otherwise $t=2$ by \cref{gra:thm:extremal}: the
baseline is $4(p+1)$, with at most four exceptional directions, each
contributing $2(r-2)$ above the baseline. A direction of depth five forces
$p\equiv1\pmod{10}$; its squared projections are a coset of $\gramu_5$
inside the nonzero squares, which also rules out a zero step. Thus zero,
one or four depth-five directions allow at most $16$, $18$ or $24$ above
the baseline, and $24$ requires four depth-five directions.

Take two depth-five directions as coordinate covectors. Squaring is a
bijection on $\gramu_5$, so after independent row scalings each row
consists of the five powers of a fixed primitive fifth root $z\in\F_p$, up
to independent signs. Permute and sign the columns so that the first row
is $(z^i)_{0\le i<5}$, and scale the second row so that its zeroth entry is
$1$. It is now $(\epsilon_iz^{\sigma(i)})_i$ with $\epsilon_0=1$ and
$\sigma(0)=0$. The cross term of $M_1=0$ gives
\[
 f(z)=0,\qquad f(Z)=\sum_{i=0}^4\epsilon_iZ^{(i+\sigma(i))\bmod5}.
\]
There are $4!\,2^4=384$ choices. Multiplication by $f$ on $\Z[Z]/(\Phi_5)$
has the matrix with entry $f_{(i-j)\bmod5}-f_{(4-j)\bmod5}$ at $(i,j)$,
$0\le i,j<4$, where $f_h$ is the coefficient of $Z^h$. If $f(z)=0$ in
$\F_p$, evaluation at $z$ annihilates the image of this matrix modulo $p$,
so $p$ divides its determinant. The absolute determinants of the $384$
cases form the set $\{0,1,11,16,25,31,81,101,121,131,625\}$. Thus for
$p>10$ every case with nonzero determinant has $p\in\{11,31,101,131\}$.
If the determinant vanishes, the coefficient vector of $f$ is constant,
because $\Phi_5$ is irreducible over $\Q$; its odd coefficient sum, with
$\epsilon_0=1$, forces all five signs to be positive. Then $i+\sigma(i)$
is a permutation modulo five, which leaves $\sigma(i)=ki\bmod5$ with
$k=1,2,3$. For $k=1$ the rows coincide. For $k=2,3$ a geometric sum gives
$M_2=20X^3Y$ or $20XY^3$, so only the axes are exceptional, both of depth
five, with excess $12$; for $k=2$ this is (e).

For the four exceptional primes, evaluating the $384$ cases at one
primitive fifth root $z$ per prime gives the following complete list of
surviving depth patterns, after discarding coincident rows; counts refer to
normalised cases, not equivalence classes.
\begin{center}\small\upshape
\begin{tabular}{rrl}
\toprule
$p$ & $z$ & exceptional depths, with multiplicities\\
\midrule
11 & 3 & $(5,5)$: 2;\quad $(5,5,5,5)$: 30\\
31 & 2 & $(5,5)$: 2;\quad $(3,5,5,5)$: 10\\
101 & 36 & $(5,5)$: 2;\quad $(3,3,5,5)$: 5\\
131 & 53 & $(5,5)$: 7\\
\bottomrule
\end{tabular}
\end{center}
Hence four depth-five directions occur only at $p=11$. At $p=31$ the best
excess with two or more depth-five directions is $20$, while at most one
such direction allows only $18$. At the other primes two or more
depth-five directions give excess at most $16$, and at most $12$ outside
$p\in\{11,31,101,131\}$. This proves the bounds in (a) and (b), and proves
(c) except when exactly one direction has depth five; outside $p=11,31$
that case is covered by the bound in (d), proved below.

For the orbits take each ordered independent pair of columns and both
relative signs, map the pair to the coordinate basis, take the smaller
representative of each sign pair and sort the list. The least list
obtained is constant on each orbit and identifies it, since every move is
invertible. The $30$ sharp cases at $p=11$ reduce to the representative in
(a), and the $10$ maximal cases at $p=31$ reduce to the orbit of the
stated witness. For the witness in (a) the integer coefficients of $M_1$
are $11,44,55$. In the directions $(1,t)$, $0\le t\le10$, and then
$(0,1)$, the residues of $M_2$ are
\[
 6,\ 2,\ 0,\ 0,\ 0,\ 7,\ 4,\ 5,\ 0,\ 10,\ 10,\ 10,
\]
and at each of its four zeros $(M_3,M_4,M_5)\equiv(0,0,5)$. Thus
$\graA=2(8\cdot2+4\cdot5)=72$. The steps are nonzero and distinct up to
sign modulo eleven, and \cref{gra:thm:moments} gives the layer formula.

It remains to prove (d). Let $p\equiv1\pmod{10}$, $p\ne11,31$, and suppose
that $M_1=0$ and exactly one direction has depth five. As above, after
scaling, reordering and signing the columns, its projection row is
$a=(1,z,z^2,z^3,z^4)$ with $z$ of order five. The componentwise powers
$a^j=(z^{ij})_{0\le i<5}$, $0\le j<5$, form a basis of $\F_p^5$, and
$a\cdot a^j=\sum_iz^{i(j+1)}$ is $5$ for $j=4$ and $0$ otherwise. A second
row $b=\sum_j\beta_ja^j$ of the row plane satisfies $a\cdot b=b\cdot b=0$.
The first condition gives $\beta_4=0$, and subtracting a multiple of $a$,
which does not change the row plane, gives $\beta_1=0$. Writing
$b=Ba^0+Ca^2+Da^3$, the second condition is $B^2+2CD=0$. If $C=0$, then
$B=0$ and $b$ is proportional to $a^3$. Otherwise scale $b$ to
\[
 b=\theta a^0+a^2-\tfrac12\theta^2a^3 .
\]
The rows $b=a^3$ and $b=a^2$ (the case $\theta=0$) are the cases $k=3,2$
above, with two directions of depth five; so $\theta\ne0$.

The other directions have projection rows $b+va$, $v\in\F_p$. In the
algebra with $X^5=1$ put
\[
 U=2\theta+2vX+2X^2-\theta^2X^3 ,
\]
whose evaluation at $X=z^i$ is $2(b+va)_i$. Since $\sum_iz^{ik}$ is $5$ or
$0$, the moment $M_j$ at $b+va$ is a nonzero scalar multiple of the
constant coefficient of $U^{2j}$. Let $f$ and $g$ be the constant
coefficients of $U^4$ and $U^6$, and put $\varsigma=\theta^5$. Exact
expansion and division give
\[
\begin{split}
 f&=64v^3-96\theta^3v^2-8\theta(\varsigma-24)v-56\theta^4,\qquad g\equiv3R\pmod f,\\
 R&=\theta(128+352\varsigma+367\varsigma^2)
   +\theta^3(-768-1122\varsigma+53\varsigma^2)v+(320+670\varsigma^2)v^2 .
\end{split}
\]
Thus the finite exceptional directions are the at most three roots of $f$,
and a root has depth at least four exactly when it is also a root of $R$.
The greatest common divisor of the three pairwise resultants in
$\varsigma$ of the coefficient quadratics of $R$ is $2^95^411^2$, so $R$ is
not identically zero in $v$. If two distinct roots of $f$ were roots of
$R$, then $R$ would be quadratic and divide $f$. Multiplying the remainder
of $f$ modulo $R$ by the square of the leading coefficient of $R$ gives
$-32\theta^4P(\varsigma)-32\theta Q(\varsigma)v$, where
\[
\begin{aligned}
 P&=252928+489984\varsigma+1480832\varsigma^2+78716\varsigma^3+9003\varsigma^4,\\
 Q&=-532480-191488\varsigma-4536064\varsigma^2-283272\varsigma^3
   +291464\varsigma^4+77\varsigma^5 .
\end{aligned}
\]
The absolute resultant of $P$ and $Q$ is
$2^{42}5^{24}11^{10}13^231^2\cdot761\cdot1297$, and the only one of its
prime factors that is $1\bmod10$, besides $11$ and $31$, is $761$. At
$p=761$ a direct check over all nonzero $\theta$ finds no common root of
$f$ and $R$. Consequently at most one of the three finite exceptional
directions has depth four, and the excess is at most $6+4+2+2=14$.

For attainment, the exact identities
\[
 \operatorname{Res}_v(f,R)=-1024\theta^3C_0(\theta^5),\qquad
 \operatorname{disc}_v(f)=65536\theta^3D_0(\theta^5)
\]
hold with
\[
\begin{split}
 C_0(\varsigma)={}&-354418688-419381248\varsigma-3296137216\varsigma^2
 -3222697984\varsigma^3+881034240\varsigma^4\\
 &+7486204\varsigma^5-75721176\varsigma^6+151118\varsigma^7+77\varsigma^8,\\
 D_0(\varsigma)={}&-27648+21492\varsigma-4356\varsigma^2+11\varsigma^3 .
\end{split}
\]
Here $C_0$ is nonconstant with nonzero constant term, and $C_0,D_0$ are
coprime, since their reductions modulo three have unchanged degrees and are
coprime. Choose $\theta\ne0$ in characteristic zero with $C_0(\theta^5)=0$.
Then $D_0(\theta^5)\ne0$, so $f$ has three distinct roots, and exactly one
of them is a root of $R$, by the preceding remainders. Take a Galois number field containing
$\theta$, the three roots and a primitive fifth root of unity. By the
Chebotarev density theorem infinitely many primes split completely in it,
and all of them are $1\bmod10$. Discard the finitely many primes at which
these identities, nonvanishings and distinctness fail to reduce. The
reduced finite roots give two directions of depth three and one of depth
at least four. After also excluding $11,31,101,131$, two depth-five
directions would allow excess at most $12$, whereas this pattern already
gives at least $16$ if its last finite depth were five; so that depth is
exactly four and $\graA=4(p+1)+14$. The first row has five distinct signed
classes $\pm z^i$, so the steps are nonzero and distinct up to sign, and
the towers are simple.
\end{proof}

The multiplication determinants agree with Sylvester resultants in all
$384$ cases, and independent fraction-free elimination of the
$120\times120$ reduced Laplacian of the witness at $n=1$ gives valuation
$70$; see
\grafiles{\artifact{five_step_sharpness_2026-11-27/verification.txt}}
Exhaustive searches give the maxima of $\graA$ at every prime
$11\le p\le181$, for all residue lists and for lists with nonzero
residues; a step divisible by $p$ in both coordinates counts as a zero
residue. For lists with $M_1=0$, a change of basis sends two independent
steps to $(1,0),(0,1)$; with sign representatives for the remaining three
steps, the key $(x^2,xy,y^2)$ is injective, and each sorted pair of the
remaining steps determines the key of the last. In every case the maximum
exceeds the ceiling $2(p+1)+16$ for $M_1\ne0$. The entries marked
$^{\mathrm a}$ and $^{\mathrm b}$ are \cref{gra:thm:fivestep}(a) and~(b).
\begin{center}\small\upshape
\begin{tabular}{rrr@{\qquad}rrr@{\qquad}rrr}
\toprule
$p$ & all & nonzero & $p$ & all & nonzero & $p$ & all & nonzero\\
\midrule
$11^{\mathrm a}$ & 72 & 72 & 61 & 260 & 260 & 127 & 520 & 520\\
13 & 64 & 58 & 67 & 280 & 280 & 131 & 540 & 540\\
17 & 88 & 80 & 71 & 300 & 300 & 137 & 568 & 562\\
19 & 88 & 84 & 73 & 312 & 306 & 139 & 568 & 568\\
23 & 96 & 96 & 79 & 328 & 328 & 149 & 608 & 608\\
29 & 128 & 128 & 83 & 344 & 344 & 151 & 620 & 620\\
$31^{\mathrm b}$ & 148 & 148 & 89 & 376 & 368 & 157 & 640 & 640\\
37 & 160 & 160 & 97 & 408 & 402 & 163 & 664 & 664\\
41 & 184 & 180 & 101 & 424 & 424 & 167 & 680 & 680\\
43 & 184 & 180 & 103 & 424 & 424 & 173 & 708 & 708\\
47 & 200 & 200 & 107 & 440 & 440 & 179 & 728 & 728\\
53 & 224 & 224 & 109 & 448 & 448 & 181 & 740 & 740\\
59 & 248 & 248 & 113 & 472 & 466 & & & \\
\bottomrule
\end{tabular}
\end{center}
The $22$ primes up to $101$ are recorded in
\grafiles{\artifact{five_step_graph_towers_2026-11-26/maxima.csv}}
At $p=131$ the maximum is $540=4(p+1)+12$, with or without zero residues;
the witness $((1,0),(0,1),(0,42),(1,20),(28,18))$ has distinct nonzero
signed classes, so $\max_Sv_{131}(\kappa_n(S))=540(131^n-1)/130-2n$, also
among simple towers:
\grafiles{\artifact{five_step_sharpness_2026-11-27/benchmark_131.txt}}
At each of the other fifteen primes from $103$ to $181$ the maximum over
lists with nonzero residues equals the maximum over simple towers. These
values are finite verifications; simple witnesses and normalisation counts
are in
\grafiles{\artifact{five_step_refinement_2026-11-27/small_benchmarks.txt}}
The polynomial identities in the proof of (d), the modular test at
$p=761$ and the witnesses below are checked in
\grafiles{\artifact{five_step_refinement_2026-11-27/symbolic.txt}}

\begin{proposition}[Exact maxima at six larger primes; \stU, finite verification]\label[proposition]{gra:prop:fivelarge}
For five steps the maxima of $\graA(S)$ over all residue lists and over
lists with nonzero residues are:
\begin{center}\small\upshape
\begin{tabular}{rrrrr}
\toprule
$p$ & $z$ & $\theta$ & all residues & nonzero residues, also simple\\
\midrule
761  & --  & --   & 3064 & 3060\\
1031 & 264 & 381  & 4142 & 4142\\
1511 & 534 & 593  & 6062 & 6062\\
1571 & 382 & 213  & 6302 & 6302\\
2281 & 342 & 170  & 9144 & 9142\\
2351 & 207 & 1124 & 9422 & 9422\\
\bottomrule
\end{tabular}
\end{center}
For the five rows with parameters, the row plane spanned by
$a=(z^i)_{0\le i<5}$ and $b=\theta a^0+a^2-\frac12\theta^2a^3$ is a simple
witness for the value in the last column. For example, at $p=1031$ the
exact maximal valuation is $4142(1031^n-1)/1030-2n$ for every $n\ge1$.
\end{proposition}

\begin{proof}
If the quartic $M_2$ has four directions of depth at least four, the row
plane meets the curve
$C_5:\sum_ix_i^2=\sum_ix_i^4=\sum_ix_i^6=0$ in $\graP^4$ in four points,
and conversely its intersections with $C_5$ are exactly the directions of
depth at least four. To enumerate $C_5(\F_p)$, use a permutation to put a
nonzero coordinate first and scale it to $1$. For squared second and third
coordinates $\alpha,\beta$, the squared last two coordinates have sum
$\sigma=-1-\alpha-\beta$ and product
$\tau=(\sigma^2+1+\alpha^2+\beta^2)/2$. The third power sum vanishes
exactly when $\sigma^3-3\sigma\tau+1+\alpha^3+\beta^3=0$; test this
identity and then whether both roots of $T^2-\sigma T+\tau$ are squares
(zero allowed). This exhausts the points after restoring permutations and
signs. For each first point up to those symmetries, pair it with every
distinct orthogonal second point; the span is isotropic, and evaluating
its remaining projective directions gives its full count of points of
$C_5$.

At $1031,1511,1571,2351$ no such plane has four points of $C_5$, so the
case without a depth-five direction gives excess at most $14$; by
\cref{gra:thm:fivestep}(d) and the certificate for two or more depth-five
directions the other cases give at most $14$, and the stated witness
attains $14$ by a direct moment check. At $2281$ every plane with four
points of $C_5$ has a zero step, so the same holds when all residues are
nonzero. At $761$ every plane with nonzero steps has at most two points of
$C_5$, so a plane without a depth-five direction has excess at most $12$.
With exactly one depth-five direction the modular check at $761$ in the
proof of \cref{gra:thm:fivestep}(d) excludes any other direction of depth
four, again giving at most $12$, and two or more give at most $12$ by the
certificate. The simple family of \cref{gra:thm:fivestep}(e) attains $12$.
At $761$ and $2281$, which are $1\bmod8$, appending a zero step to an
optimal four-step list (\cref{gra:thm:four}) attains excess $16$, the
general ceiling of \cref{gra:thm:fivestep}(c).
\end{proof}

The search counts, which are for inspected pairs of points rather than
plane orbits, are in
\grafiles{\artifact{five_step_refinement_2026-11-27/depth_four.txt}}

\paragraph{Seven steps and prime step counts.}
For $m=7$ the bound \eqref{gra:eq:tbound} is $6(p+1)+48$.

\begin{theorem}[Seven steps; \stU]\label{gra:thm:seven}
Let $m=7$ and $p>14$. Then $\graA(S)\le6(p+1)+38$, and
$\graA(S)\le6(p+1)+36$ if $p\not\equiv1\pmod{14}$. If $M_1=M_2=0$ and two
directions have depth seven, these are the only directions of depth above
three, and $\graA(S)=6(p+1)+16$. This case occurs in simple graphs for
every $p\equiv1\pmod{14}$: if $z$ has order seven modulo $p$, take the
steps $(z^i,z^{3i})$, $0\le i<7$.
\end{theorem}

\begin{proof}
For $t\le2$, \eqref{gra:eq:tbound} gives at most
$4(p+1)+40<6(p+1)+36$. When $t=3$ there are at most six exceptional
directions. If none has depth seven, each adds at most six above the
depth-three baseline. If exactly one has depth seven, it adds eight and
the other five add at most six each. A depth-seven direction requires
$14\mid p-1$ by \cref{gra:thm:extremal}. It remains to treat two or more
depth-seven directions with $t=3$.

Take two depth-seven directions as coordinate axes. Each projection row
has squared entries in a coset of $\gramu_7$, and squaring is bijective on
$\gramu_7$. Row scalings, column permutations and column signs put the
rows in the form $a_i=z^i$, $b_i=\epsilon_iz^{\sigma(i)}$, where $z$ has
order seven, $\sigma(0)=0$ and $\epsilon_0=1$. The coefficient of
$X^2Y^2$ in $M_2$ requires $f(z^2)=0$ for
$f(Z)=\sum_{i=0}^6Z^{(i+\sigma(i))\bmod7}$. The absolute norms of $f$ over
$\Q(\zeta_7)$ for the $6!=720$ permutations, computed as determinants of
multiplication by $f$ on $\Z[Z]/(\Phi_7)$ (entry
$f_{(i-j)\bmod7}-f_{(6-j)\bmod7}$ at $(i,j)$, $0\le i,j<6$) and
independently as exact products of conjugates, are:
\begin{center}\small\upshape
\begin{tabular}{lrrrrrrrr}
\toprule
absolute norm & 0 & 49 & 343 & 392 & 2744 & 3479 & 8281 & 117649\\
multiplicity & 19 & 441 & 42 & 126 & 28 & 42 & 21 & 1\\
\bottomrule
\end{tabular}
\end{center}
Evaluation at $z^2$ annihilates the image modulo $p$, so $p$ divides the
norm. The only prime $p>14$ dividing a nonzero entry is $71$. Elsewhere
the norm is zero, so irreducibility of $\Phi_7$ forces all seven
coefficients of $f$ to equal one. The cross moment $\sum_ia_ib_i$ is then a
signed sum of the seven distinct seventh roots of unity. Among sign
polynomials with constant coefficient $+1$ the absolute norms are
$0,64,512$, with multiplicities $1,49,14$; the only one vanishing at an
odd prime has every sign positive. The coefficients of $X^3Y$ and $XY^3$
in $M_2$ give the two further exponent lists $3i+\sigma(i)$ and
$i+3\sigma(i)$ modulo seven, and the same norm table applies after
relabelling. Hence all three exponent lists are permutations, which leaves
$\sigma(i)=ki$ with $k\in\{1,3,5\}$; the identity does not span. For
$k=3,5$ geometric sums give $M_3=105X^2Y^4$ or $105X^4Y^2$, respectively.
As $p>14$ this has only the two axis roots, both of depth seven, and
$\graA=6(p+1)+16$.

At $p=71$ take $z=20$ and test all $720\cdot64$ permutation and sign
choices on the mixed coefficients of $M_1$ and $M_2$. After dependent rows
are excluded, precisely $k=3,5$ with positive signs survive. This
completes the classification. The steps $(z^i,z^{3i})$ span and have
distinct nonzero signed classes, so they give simple graphs at every
level, and their depth pattern proves the attainment.
\end{proof}

The norm and moment certificate is
\grafiles{\artifact{seven_step_graph_towers_2026-11-27/certificate_output.txt}}
Exhaustive searches give the following maxima of $\graA$ for seven steps.
At $p=29$ the search covers all lists with $M_1=M_2=0$: there are twelve
normalised candidates, the best has $\graA=196=6(p+1)+16$, and lists with
$t\le2$ have $\graA\le160$; a witness is the basis together with
$(2,3),(6,2),(10,24),(11,21),(12,10)$, whose only exceptional directions
are $(1,7)$ and $(1,24)$, both of depth seven. At $p=17,19,23,31$ the
search covers all lists with $M_1=0$; there the depth is at most six and
lists with $M_1\ne0$ have $\graA\le2(p+1)+20$, below the best value in
each class.
\begin{center}\small\upshape
\begin{tabular}{rrrr}
\toprule
$p$ & all residues & nonzero residues & simple graphs\\
\midrule
17 & 88 & 82 & 82\\
19 & 120 & 94 & 94\\
23 & 104 & 104 & 104\\
29 & 196 & 196 & 196\\
31 & 192 & 140 & 140\\
\bottomrule
\end{tabular}
\end{center}
Each column gives exact maxima at every level within its class by
\cref{gra:thm:moments}; for example the simple maximum at $p=29$ is
$196(29^n-1)/28-2n$. The search logs with explicit simple witnesses are
\grafiles{\artifact{seven_step_graph_towers_2026-11-27/search_output.txt}\\
\artifact{seven_step_graph_towers_2026-11-27/search_lower_output.txt}}

\begin{theorem}[Prime step counts; \stU]\label{gra:thm:primestep}
Let $m\ge7$ be prime, $k=(m-1)/2$ and $p>\max(m^{m-1},27^k)$. If $M_1=0$,
there are at most two directions of depth $m$. Consequently
\[
 \graA(S)\le2k(p+1+2m-2k)-2(m-3),
\]
which is $2(m-3)$ below the outer bound of \eqref{gra:eq:tbound}.
\end{theorem}

\begin{proof}
Suppose there are three distinct depth-$m$ directions. Choose a primitive
$m$th root of unity $z\in\F_p$ and, as above, normalise the projection rows
to $a_i=z^i$, $b_i=\epsilon_iz^{\sigma(i)}$ and $c_i=\eta_iz^{\tau(i)}$.
Lift by replacing $z$ with a complex primitive $m$th root of unity
$\zeta$. The cross moment $\sum_ia_ib_i$ and all $3\times3$ minors of the
three-row matrix have zero reduction at the prime $(p,\zeta-z)$. If one of
them is nonzero, its integral norm is divisible by $p$. Each conjugate of
the cross moment has absolute value at most $m$, and each conjugate of a
minor at most $3\sqrt3$ by Hadamard's inequality, so the norms are at most
$m^{m-1}$ and $27^k$, both less than $p$. Hence the cross moment and the
minors vanish exactly.

Reducing exponents modulo $m$, the cross moment is a sum of $m$ signs on
powers of $\zeta$. Its coefficient vector is constant by irreducibility of
$\Phi_m$. The coefficient sum is odd with absolute value at most $m$, so
the constant is $\pm1$ and all signs agree; absorb the common sign in $b$.
The three rows are pairwise independent, as their reductions modulo $p$
already show, and the vanishing minors give $c=\alpha a+\beta b$ with
$\alpha\beta\ne0$. Every ratio $w=b_i/a_i$ satisfies
$\abs w=\abs{\alpha+\beta w}=1$. Expanding the second condition on the
unit circle gives a real affine line in $w$, which meets the unit circle
in at most two points. Thus $\sigma$ acts by at most two translations
$u,v$ modulo $m$. If it translates a set $I$ by $u$ and the complement by
$v$, bijectivity gives $I+u=I+v$; a proper nonempty $I$ is not invariant
under a nonzero translation modulo a prime. So only one translation
occurs and $a,b$ are proportional, a contradiction.

For $t=k$ there are at most $2k$ exceptional directions, and all but two
have depth at most $m-1$; relative to the outer bound of
\eqref{gra:eq:tbound} this saves at least $2(2k-2)=2(m-3)$. If $t<k$, the
bound \eqref{gra:eq:tbound} at $t=k-1$ already saves $2(p+5)$.
\end{proof}

\begin{remark}[The cases $m=5$ and $m=7$]
The proof applies verbatim to $m=5$, for $p>729$, and gives
$\graA(S)\le4(p+1)+20$ and at most two depth-five directions when $M_1=0$.
Both are superseded: \cref{gra:thm:fivestep}(c),(d) give
$\graA(S)\le4(p+1)+16$ for every prime $p>10$ other than $11$ and $31$, and
its proof gives at most two depth-five directions whenever $M_1=0$ and
$p\notin\{11,31,101,131\}$. For $m=7$, where the theorem needs
$p>7^6=117649$, its bound $6(p+1)+40$ is superseded by the bound $6(p+1)+38$
of \cref{gra:thm:seven} for every $p>14$. Its structural assertion at $m=7$,
at most two directions of depth seven when $M_1=0$, is not covered by
\cref{gra:thm:seven}, whose classification of two depth-seven directions
assumes $t=3$, that is, $M_2=0$ as well. As a bound, \cref{gra:thm:primestep}
is new for $m\ge11$.
\end{remark}

Thus for odd $m$ the outer bound of \eqref{gra:eq:tbound} is attained for
$m=3$ exactly when $p\equiv1\pmod6$, for $m=5$ only at $p=11$
(\cref{gra:thm:fivestep}), never for $m=7$ (\cref{gra:thm:seven}), and never
for prime $m\ge11$ once $p>\max(m^{m-1},27^k)$.

\subsection{How late can the eventual law start?}\label{gra:sec:onset}
\sources{Effective_Graph_Towers; grh_method_transfers; spanning-tree-sharp-cutoff}

The \emph{onset} of a tower is the least level from which the eventual
polynomial agrees with $\nu_n$. This subsection shows that the onset is
unbounded for a fixed prime and a fixed number of loops, that even the
leading coefficient of the eventual polynomial is not a continuous function
of the voltages, and that in rank two the onset is nevertheless at most
logarithmic in the size of the voltages; the worst case grows linearly in
the bit length exactly when the number of loops is at least $\max(3,p)$.
We write $C_k$ for a cyclic group of order $k$.

\begin{theorem}[The first stable level can be prescribed; \stU]\label{gra:thm:delayed}
Fix a prime $p\equiv3\pmod4$ and an integer $k\ge1$. On
$(\Z/p^n\Z)^2$ take one step in each of the directions $p^ke_1$ and $e_2$,
and $p$ copies of the step $e_1$. Put $s=\min(k,n)$ and $C=2(p+1)/(p-1)$.
At every level $n\ge1$,
\[
 \nu_n=2sp^n+C(p^n-p^s)+(p^s-1)+2s-2n .
\]
The eventual polynomial is
\[
 P_k(X,Y)=(2k+C)X-2Y+(1-C)p^k+2k-1 .
\]
It gives $\nu_n=P_k(p^n,n)$ from level $k$ onward. For $k\ge2$ it fails at
level $k-1$, so $k$ is the exact first stable level.
\end{theorem}

These connected towers all have the same number $p+2$ of base loops for
fixed $p$; in particular five loops suffice at $p=3$ to prescribe any
positive starting level.

\begin{proof}
Since $p\equiv3\pmod4$, \cref{gra:eq:aniso} gives
$v_p(\grah(U)+\grah(V))=2/e_j\le1$ for every pair $(U,V)$ of $p$-power
roots of unity of joint conductor $p^j>1$. The eigenvalue at a character $(X,Y)$ is
$\grah(X^{p^k})+\grah(Y)+p\,\grah(X)$. The final term has valuation
strictly above one, with the convention $v_p(0)=+\infty$. Thus whenever
$(X^{p^k},Y)\ne(1,1)$, the joint conductor of this image alone determines
the valuation.

The map $(X,Y)\mapsto(X^{p^k},Y)$ has fibres of size $p^s$ and image
isomorphic to $C_{p^{n-s}}\times C_{p^n}$. For $1\le j\le n-s$ the exact
conductor-$p^j$ shell of the image has $p^{2j}-p^{2j-2}$ elements; their
aggregate valuation is therefore $2(p+1)p^{j-1}$. For $n-s<j\le n$ there
are $p^{n-s}e_j$ elements in the shell, contributing $2p^{n-s}$. Summing
these shell contributions and multiplying by $p^s$ gives
$C(p^n-p^s)+2sp^n$. On the kernel with the trivial character removed, the
eigenvalue is $p\,\grah(X)$ for each nontrivial $p^s$-th root $X$. There
are $p^s-1$ such roots, and $\prod_{X\ne1}(1-X)=p^s$ gives
$\sum_{X\ne1}v_p(\grah(X))=2s$. The matrix-tree denominator contributes
$-2n$. This proves the formula.

For $n\le k$ it reads $\nu_n=(2n+1)p^n-1$, and for $n\ge k$ it equals
$P_k(p^n,n)$. The discrepancy one level before $k$ is
\[
 \nu_{k-1}-P_k(p^{k-1},k-1)=(p+1)p^{k-1}-2>0 .
\]
Distinct polynomials in $p^n,n$ cannot agree on an infinite tail: the term
of highest degree in $p^n$ dominates all lower exponential orders, even
with polynomial factors in $n$. This establishes the claimed onset for the
unique eventual polynomial.
\end{proof}

\begin{remark}\label[remark]{gra:rem:delayed}
(a) The growing onset is visible directly in the formal specialisations.
For primitive $(a,b)\in\Z_p^2$, substitute $X=(1+z)^a$ and $Y=(1+z)^b$ in
the Laplacian polynomial and reduce modulo $p$. The first nonzero degree
is $2p^{\min(k+v_p(a),v_p(b))}$: the two summands that survive reduction
have these respective initial orders, and in an equal-order tie the same
anisotropic sum of squares prevents cancellation. The maximum is exactly
$2p^k$, attained at $(1,0)$. Thus even within this family with no
identically zero primitive specialisation, neither the maximal degree
required to see a unit nor the onset admits a bound in terms of $p$, the
dimension and the base edge count alone. Replacing $p^ke_1$ by the zero
step gives the law $(2n+1)p^n-1$ at every level, despite agreement of the
voltage data modulo $p^k$.

(b) For $p=3$ and $k=3$ the stable law is $10\cdot3^n-2n-76$ and the first
five values are $8,44,188,726,2344$
(\artifact{Effective_Graph_Towers/delayed_tower_results.txt}).

(c) There is a version with distinct steps. For $p\ge7$ in the same
congruence class, replace the $p$ copies of $e_1$ by $(1+ip^k)e_1$ for
$0\le i<p$. Every valuation stays the same. Indeed, writing $Z=X^{p^k}$,
the replacement contributes $2p-X\Phi_p(Z)-X^{-1}\Phi_p(Z^{-1})$. At $Z=1$
this is $p\,\grah(X)$. If $Z$ has order $p$ the two cyclotomic factors
vanish. If its order is larger, each has valuation $(p-1)v_p(Z-1)$ by
$\Phi_p(Z)=(Z^p-1)/(Z-1)$. Hence for $Z\ne1$ the new contribution has
valuation strictly greater than $2v_p(Z-1)$, which in turn is at least
$v_p(\grah(Z)+\grah(Y))$. For $n>k$ all steps are nonzero and distinct up
to sign, so the graphs at those levels are simple.
\end{remark}

\begin{theorem}[The leading coefficient can change after an arbitrary level; \stU]\label{gra:thm:leading}
Fix $p=3$. Take three copies of each of $(1,0)$ and $(0,1)$ and one
additional step $(t,0)$, and write $\nu_n(t)$ for the resulting valuation.
For $k\ge1$ put
\[
 B(n)=9^n+4\cdot3^n-2n-5,\qquad
 D_k(n)=(2(n-k)+1)3^{n+k}+(2(n-k)+4)3^k-2n-5 .
\]
Then $\nu_n(0)=B(n)$ for every $n\ge1$, and
\[
 \nu_n(3^k)=\begin{cases}B(n),&1\le n\le k,\\ D_k(n),&n\ge k.\end{cases}
\]
Both expressions agree at $n=k$. The least positive starting level for the
eventual formula $D_k$ is exactly $k$.
\end{theorem}

\begin{proof}
Index a character by its two coordinate roots of unity, and let $a,b$ be
their conductor exponents, with exponent $0$ for the trivial root. For a
nontrivial character write $h=\max(a,b)$ and express the roots as
$z^u,z^v$ for $z$ primitive of order $3^h$. By \cref{gra:eq:aniso} at
$p=3$, the grid symbol $F=\grah(z^u)+\grah(z^v)$ has
$v_3(F)=2/\varphi(3^h)=3^{1-h}$. If $a\le k$, the full symbol equals $3F$
and has valuation $1+3^{1-h}$. If $a>k$, the extra summand
$\grah(z^{3^ku})$ has valuation $3^{1-a+k}\le1$, strictly below that of
$3F$, and gives the valuation of the sum.

For $n\le k$ there are $9^n-1$ nontrivial characters, and the
$8\cdot9^{h-1}$ characters at conductor $h$ contribute $8\cdot3^{h-1}$
from the grid valuation. Summing and subtracting the matrix-tree
correction $2n$ gives $B(n)$. This argument holds for every $n$ at $t=0$.
For $n\ge k$, the rectangle $a\le k$ contains $3^{n+k}-1$ nontrivial
characters; their unit contribution is this count. Their grid valuations
sum to $4(3^k-1)$ for $h\le k$ and to $2(n-k)3^k$ for $h>k$, since then
$b=h$. Outside the rectangle, the $2\cdot3^{a-1+n}$ characters with given
$a>k$ contribute $2\cdot3^{k+n}$. Adding these sums and subtracting $2n$
gives $D_k(n)$. For $k\ge2$ the difference at the preceding level is
\[
 B(k-1)-D_k(k-1)=4\cdot3^{2k-2}-2\cdot3^{k-1}>0 .
\]
This proves minimality, including $k=1$ by the positive-level convention.
\end{proof}

For $k=3$ the successive valuations at levels $1$ through $5$ are
$14,108,826,6710,33006$
(\artifact{Effective_Graph_Towers/leading_results.txt}).

\begin{corollary}[No digit-oracle algorithm for the leading coefficient; \stU]\label[corollary]{gra:cor:oracle}
Write $X=3^n$, and denote the coefficient of $X^2$ in the eventual
polynomial by $\mu$ and that of $nX$ by $\lambda$. In the family of
\cref{gra:thm:leading} the pair $(\mu,\lambda)$ is $(1,0)$ at $t=0$ and
$(0,2\cdot3^k)$ at $t=3^k$. Thus the degree in $X$ drops from two to one,
$\lambda$ is unbounded along a sequence for which all voltage digits
stabilise, and $\mu$ is discontinuous at $0$ on the space of $3$-adic
inputs. No algorithm that queries the $3$-adic digits of the voltages
computes $\mu$ on all inputs, and the same holds for finite programs
computing the digits, even with a promise that their values are ordinary
integers.
\end{corollary}

\begin{proof}
Any digit-query algorithm halting at $0$ has read only finitely many
digits. A sufficiently large $3^k$ yields the same answers to all these
queries, yet the required leading coefficient differs. For programs: given
a machine, output the residues of $3^j$ if it first halts at time
$j\ge1$, and the residues of zero if it never halts. Computing the residue
modulo $3^r$ needs at most $r$ simulation steps. An algorithm for $\mu$ on
such programs would decide halting. A finite integer encoding, as in
\cref{gra:thm:effective}, distinguishes these inputs.
\end{proof}

\begin{remark}[A fixed simple base graph]\label[remark]{gra:rem:simplebase}
The formulas of \cref{gra:thm:leading} hold with a simple graph at every
level. Replace each of the seven base loops by a cycle of length four,
using separate internal vertices for each cycle; put the original voltage
on one oriented edge of that cycle and zero on the other three. This
replaces every edge in each lifted graph by a path of length four. For a
connected multigraph with $E$ edges and $V$ vertices, the resulting
spanning-tree count is $4^{E-V+1}$ times the original count: the intact
paths project to a spanning tree, and on each remaining path one of the
four edges must be removed. The argument also applies to subdivided loops.
The factor is a $3$-adic unit, so the least onset $k$, the discontinuity
of $\mu$ and the unbounded values of $\lambda$ all persist on covers of
this single simple base.
\end{remark}

We turn to upper bounds in rank two.

\begin{theorem}[Quantitative shell theorem; \stU]\label{gra:thm:shell}
Let $F\ne0$ be an integral Laurent polynomial in two variables, supported
in $[-H,H]^2$ with $H\ge1$. Assume that $F$ is nonzero at all nontrivial
pairs of $p$-power roots of unity, and let $S_h$ be the sum of
$v_p(F(x,y))$ over the pairs of joint order exactly $p^h$. Set
\[
 K=\ceil{\log_p(8H^2+1)},\qquad M=\ceil{\log_p(2H+1)},\qquad h_0=K+M+2 .
\]
There exist rational $s_2,s_1,t_1,s_0$ such that for every $h\ge h_0$
\[
 S_h=s_2p^{2h}+(s_1+t_1h)p^h+s_0 .
\]
The constants are independent of the level, and the bound $h_0$ depends
only on the exponents and $p$.
\end{theorem}

\begin{proof}
Dividing $F$ by the common $p$-power $p^g$ of its coefficients adds
$g(p^{2h}-p^{2h-2})$ back to the shell. Collect equal exponents and assume
henceforth that $F=\sum_ic_iX^{a_i}Y^{b_i}$ has a unit coefficient. Fix
compatible primitive roots $\zeta_h$. Each orbit of the unit action, of
size $e_h$, contributes $\ord_h(F)=e_hv_p(F)$. Each orbit is uniquely
represented in one of two charts: $(\zeta_h,\zeta_h^z)$ for any $z$ modulo
$p^h$, or $(\zeta_h^z,\zeta_h)$ with $p\mid z$. In the second chart
interchange the two exponent coordinates below.

The projected exponents in the first chart are $a_i+b_iz$. Take as
\emph{centres} the distinct $p$-integral rational solutions of
$a_i+b_iz=a_j+b_jz$ lying in the chart. Write each centre as a reduced
fraction $A/B$ with $p\nmid B$ and $\abs A,\abs B\le2H$. The difference of
two distinct centres has nonzero numerator of absolute value at most
$8H^2$. Hence their disks $z\equiv A/B\pmod{p^K}$ are disjoint. In the
restricted chart each disk is entirely inside or outside the chart, since
$K\ge1$.

Outside these disks the projected exponents are pairwise distinct modulo
$p^J$, where $J=K+M$. In fact a nonzero slope difference $b_i-b_j$ has
valuation $<M$, and the distance factor $z-(a_j-a_i)/(b_i-b_j)$ has
valuation $<K$; this covers centres outside the chart too. A nonzero
constant difference has valuation $<M$. For each fixed $z\bmod p^J$ in this
complement let $n_i\in[0,p^J)$ be the distinct exponent residues. The
polynomial $\sum_ic_i(1+T)^{n_i}$ is nonzero modulo $p$ by the distinct
degrees and the unit coefficient. Let its first nonzero coefficient occur
at $j<p^J$. For any lift of $z$ and $h\ge h_0$, the orbit contributes
exactly $j$. Indeed, modulo $(p,T^{p^J})$ the binomial expansion of each
monomial depends only on its exponent modulo $p^J$, and one may choose
nonnegative representatives modulo $p^h$ in forming the actual polynomial.
At $T=\zeta_h-1$ a coefficient divisible by $p$ has order at least
$e_h\ge p^J$, and all later terms have larger order than $j$. Each class
contributes $jp^{h-J}$ to $S_h$.

Fix a centre $A/B$. Applying the Galois automorphism
$\zeta_h\mapsto\zeta_h^B$ expresses the same valuation as that of
\[
 P(\zeta_h,W),\qquad W=\zeta_h^{Bz-A},\qquad
 P(X,Y)=\sum_ic_iX^{Ba_i+Ab_i}Y^{b_i},
\]
up to multiplication by a monomial of valuation zero. Shift each exponent
axis by its minimum, so that $P$ is a polynomial of degrees $D\le8H^2$ in
$X$ and $E\le2H$ in $Y$. The exponent transformation is injective, so $P$
still has a unit coefficient. As $z$ ranges over the disk modulo $p^h$,
$W$ ranges exactly once over the roots of order dividing $p^{h-K}$, since
$B$ is a unit. Write $r$ for the conductor exponent of $W$, with $r=0$ for
$W=1$, and put $R=M+1$.

For $0\le r\le R$, use $e_0=1$ and write
$P(1+T,W)=\sum_{j=0}^Dd_j(W)T^j$. This polynomial is nonzero: its
vanishing would give a torsion zero of $F$ at the pair represented by
$z$. Its nonzero coefficients have valuations in $e_r^{-1}\Z$. Let
$\mu_W$ be their minimum and let $j_W$ be the first index attaining it. We
have $e_h/e_r>D$ (for $r\ge1$ the quotient is $p^{h-r}$, and for $r=0$ use
$e_h\ge p^{h-1}$), so terms with larger coefficient valuation cannot
minimise; the unique minimum has order $e_h\mu_W+j_W$. The set of all $W$
with $r\le R$ is the same for every $h\ge h_0$. Its total is a constant
multiple of $p^h$ plus a constant.

For $R<r\le h-K$, put $k=h-r$. Then $\ord_h(W-1)=p^k>D$ and $e_r>E+1$.
Expand $P(X,1+U)=\sum_{i=0}^EP_i(X)U^i$, and take the smallest $l$ such
that $P_l(X)$ is nonzero modulo $p$; such an $l$ exists. Let $d\le D$ be
the order of $P_l(1+T)$ at $T=0$ modulo $p$. Its value at $X=\zeta_h$ has
order exactly $d$. Terms with $i<l$ have order at least
$e_h=e_rp^k>lp^k+d$, and terms with $i>l$ have order at least
$(l+1)p^k>lp^k+d$. The orbit contribution is exactly $lp^k+d$ for every
root of conductor $p^r$. The contribution from all these roots is
\[
 \sum_{k=K}^{h-R-1}e_{h-k}(lp^k+d)=l(h-R-K)e_h+d(p^{h-K}-p^R).
\]
Adding the finitely many disks in each chart, their complements and the
removed content proves the shell formula.
\end{proof}

\begin{corollary}[Logarithmic bound for the onset in rank two; \stU]\label[corollary]{gra:cor:logbound}
Let $F$ be the determinant of the Laurent Laplacian of a finite connected
base graph with voltages in $\Z^2$ whose finite lifts are connected, and
let $h_0$ be as in \cref{gra:thm:shell} for the support of $F$. Then from
$n=h_0$ onward
\[
 \nu_n=\alpha p^{2n}+(\beta+\gamma n)p^n+\delta n+\epsilon
\]
for rational constants. For a bouquet whose step coordinates are at most
$2^b$ in absolute value, the onset is at most $3b\log_p2+O_p(1)$. For a
fixed base graph the same estimate holds with a base-dependent additive
constant.
\end{corollary}

\begin{proof}
Connectedness gives invertible blocks at every nontrivial character, so
\eqref{gra:eq:shells} applies with $d=2$; summing the shell formula gives
the displayed shape from $n=h_0$ onward. For a bouquet the support of the
Laurent Laplacian consists of zero and the signed voltage vectors. For a
fixed base graph the coordinates in the determinant support are bounded by
the vertex count times the largest voltage coordinate, by determinant
expansion.
\end{proof}

\begin{remark}\label[remark]{gra:rem:five}
Five consecutive values beyond the bound $h_0$ suffice to recover the law:
the evaluation matrix on $1,n,p^n,np^n,p^{2n}$ at five consecutive levels
is nonsingular, since these sequences solve the order-five recurrence with
characteristic polynomial $(x-1)^2(x-p)^2(x-p^2)$ and are linearly
independent. Checking the finitely many preceding levels then gives the
least onset. This is a bound on conductor precision; enumerating all the
characters there can still be expensive.
\end{remark}

\begin{remark}[The coefficient of $p^{2n}$]\label[remark]{gra:rem:content}
The coefficient $\alpha$ of $p^{2n}$ in \cref{gra:cor:logbound} is exactly
the content valuation $g=\min_iv_p(c_i)$ of the collected Laurent
determinant: restoring the removed power contributes exactly $g(p^{2n}-1)$
to the cumulative sum, and every other summand grows as $O(np^n)$. This
extends the discontinuity of \cref{gra:cor:oracle} to every prime. On a
bouquet of $2p+1$ loops put $p$ copies each of $e_1,e_2$ and one step
$p^se_1$ with $s\ge1$, and compare with the same assignment where the
final step is zero. Both towers are connected, and the graphs agree
through level $s$. The zero-step determinant is
$p(\grah(X)+\grah(Y))$ and has content valuation one. The extra summand
$\grah(X^{p^s})$ changes the content valuation to zero, because its two
nonconstant exponents are new. So the two eventual laws have different
$p^{2n}$ coefficients. Moreover, once $s-4$ is at least the bound $h_0$ of
the zero-step tower, the zero-step law is exact throughout $[s-4,s]$; if
the other law were exact from $s-4$ onward the two would agree on these
five consecutive levels, which \cref{gra:rem:five} excludes. Hence the
onset of the perturbed tower is at least $s-3$. At $p=3$ this is the
family of \cref{gra:thm:leading}, where the onset is exactly $s$.
\end{remark}

\begin{theorem}[Classification by the number of loops; \stU]\label{gra:thm:loops}
Let $L_{p,m}(b)$ be the largest onset over connected rank-two towers on one
bouquet of $m\ge2$ loops with step coordinates bounded by $2^b$ in
absolute value; repeated directions and zero voltages are allowed. As
$b\to\infty$,
\[
 L_{p,m}(b)=\begin{cases}O_p(1),&2\le m<\max(3,p),\\
 \Theta_p(b),&m\ge\max(3,p).\end{cases}
\]
The constants in the upper bound can be taken independent of $m$. For odd
$p$ and $m<p$ the onset is at most one, by \cref{gra:thm:cutoff}.
\end{theorem}

\begin{proof}
The upper bound $O_p(b)$, with constants independent of $m$, is
\cref{gra:cor:logbound}.

\emph{The bounded range.} Let $p$ be odd and $m<p$. Connectedness forces
the steps to span $\F_p^2$, so \cref{gra:thm:cutoff} gives
$\nu_n=C(p^n-1)/(p-1)-2n$ for every $n\ge1$, and the onset is at most
one. For $p=2$, $m=2$ the two steps
form an invertible basis modulo $2$, hence modulo $2^n$ for every $n$. A
change of coordinates gives the same square grid at every fixed level for
all choices of voltages, so the support bound $h_0$ evaluated at $H=1$
suffices.

\emph{A test for the coefficient of $np^n$.} For the linear range it is
enough to work at the threshold value $m=\max(3,p)$ and pad with zero
voltages. If a Laurent polynomial $F$ as in \cref{gra:thm:shell} is
primitive modulo $p$, the coefficient of $np^n$ in the cumulative formula
is positive precisely when some nonzero rational direction $(u,v)$, taken
as a primitive integer representative, satisfies
\begin{equation}\label{gra:eq:test}
 F(Z^u,Z^v)\equiv0\pmod p\quad\text{as a Laurent polynomial in }Z .
\end{equation}
Indeed, the coefficient of $hp^h$ in $S_h$ is the sum of the nonnegative
disk contributions $l(p-1)/p$. A disk contributes positively exactly when
$P(X,1)$ vanishes modulo $p$, which is \eqref{gra:eq:test} after the
monomial shift and the change of coordinates defining the disk. Conversely
any identity \eqref{gra:eq:test} must group at least two distinct
exponents with the same projection, because there is a unit coefficient;
this yields one of the rational centres, using the chart where the
direction has a unit pivot. Summing $hp^h$ gives a positive multiple of
$np^n$ plus a multiple of $p^n$ and a constant.

\emph{The case $m=p$ odd.} Use $p-2$ copies of $(1,0)$, one copy of
$(1,1)$ and one copy of $(1+\delta,2)$, first with $\delta=0$ and then
with $\delta=p^s>4$. In both cases the content is a unit and the first two
directions ensure connectedness. In the first case projection along
$(u,v)=(1,0)$ makes all $p$ exponents equal to one, so the test is
positive. For the second case a vanishing projection must have all $p$
absolute projected exponents equal to one another and nonzero: an occupied
positive absolute frequency $f$ contributes the negative of its
multiplicity to the coefficient of $Z^f$, this multiplicity must equal
$p$, and a nonzero direction cannot annihilate every step. Thus $u\ne0$
and $\abs{u+v}=\abs u$, so $v=0$ or $v=-2u$. The last projected absolute
value is then $\abs{(1+\delta)u}$ or $\abs{(\delta-3)u}$, both different
from $\abs u$. The $np^n$ term vanishes after perturbation.

\emph{The case $p=2$, $m=3$.} Use the three steps
$(1,0),(1,1),(\delta,1+2\delta)$ with $\delta=0$ or $2^s$. Again the
content is a unit. In characteristic two every positive absolute projected
frequency must occur an even number of times. Three steps spanning the
plane therefore give a zero Laurent polynomial exactly when one projected
value is zero and the other two have equal nonzero absolute value. A
direction annihilating one step projects the other two in proportion to
their determinants with that step. For $\delta=0$ the three pairwise
absolute determinants all equal one, and such a direction exists. For
$\delta>0$ the determinants are $1,1+\delta,1+2\delta$, all distinct. Thus
this pair also has different $np^n$ coefficients.

\emph{Conclusion.} In both constructions the perturbed tower agrees
exactly with the fixed initial tower up through level $s$. Once the fixed
initial tower has stabilised by $s-4$, the argument of
\cref{gra:rem:five} forces onset at least $s-3$ for the perturbed tower.
The heights are at most $1+2p^s$, allowing $s=b\log_p2+O_p(1)$. Together
with the upper bound this proves the classification.
\end{proof}

Sample shell sums for \cref{gra:thm:shell} and the vanishing directions
used in the two constructions are recorded in
\artifact{grh_method_transfers/checks/onset_output.txt}.

\subsection{Stabilisation of the prime-to-$p$ part}\label{gra:sec:units}
\sources{spanning-tree-unit-stabilization}

For a nonzero rational $z$ put $\graunit(z)=z/p^{v_p(z)}$. The valuations
$\nu_n$ describe the $p$-part of the spanning-tree counts; the following
theorem describes the remaining factor, also for nonabelian deck groups.
Let $\Gamma$ be a compact $p$-adic analytic pro-$p$ group and $W$ an open
normal uniform subgroup. Let $X$ be a finite connected regular graph and
$(X_n)_{n\ge0}$ a tower of connected Galois covers of $X$, compatible
under the quotient maps, with deck groups $\Gamma/W^{p^n}$. Write
$\kappa_n$ for the number of spanning trees of $X_n$ and
$U_n=\graunit(\kappa_n)$.

\begin{theorem}[Sharp stabilisation of the prime-to-$p$ part; \stU]\label{gra:thm:units}
For every $n\ge1$,
\[
 U_n\equiv U_{n-1}\pmod{p^n}\quad(p\text{ odd}),\qquad
 U_n\equiv U_{n-1}\pmod{2^{n+1}}\quad(p=2,\ n\ge2),
\]
and $U_1\equiv U_0\pmod2$ for $p=2$. Consequently $U_n$ converges in
$\Z_p^\times$, and $U_N$ determines the limit modulo $p^{N+1}$ for odd $p$
and modulo $2^{N+2}$ for $p=2$ ($N\ge1$). These moduli are sharp at every
level, even for towers of simple graphs (\cref{gra:prop:unitsharp}).
\end{theorem}

\begin{proof}
\emph{Reduction to one step.} Let $G=\Gamma/W^{p^n}$ and
$K=W^{p^{n-1}}/W^{p^n}$, so that $X_n\to X_{n-1}$ is the quotient by $K$.
The kernel $K$ is elementary abelian and central in $W/W^{p^n}$, and every
element of $K$ is the $p^{n-1}$st power of an element of $W/W^{p^n}$, by
standard properties of uniform pro-$p$ groups. Only these properties are
used: for each nontrivial character of $K$ it suffices to have one such
power that the character detects, with the root centralising $K$ and of
order dividing $p^n$.

Choose $L=\Q(\zeta_{p^M})$ with $M\ge n$ and $p^M$ a multiple of the
exponent of $G$. Let $A$ be its Galois group over $\Q$, identified with
the local Galois group at $p$ and with $(\Z/p^M\Z)^\times$. Decompose the
lifted Laplacian into blocks using the regular representation of $G$ on
each fibre over the base. For each irreducible representation $\rho$ the
block $D_\rho$ appears $d_\rho=\dim\rho$ times; it is invertible for
nontrivial $\rho$, because connectivity makes the constants the whole
kernel. The matrix-tree theorem cancels the factors from representations
trivial on $K$ and gives
\[
 \frac{\kappa_n}{\kappa_{n-1}}=\frac1{\abs K}
 \prod_{\rho:\,K\not\subseteq\ker\rho}(\det D_\rho)^{d_\rho}.
\]
All blocks are Hermitian in unitary bases, and their determinants are
real and positive. The Galois action on these determinants agrees with its
action on the representations, since the voltage Laplacian has rational
group-ring coefficients.

\emph{The stabiliser.} Fix $\rho$ in this product, let $H\le A$ be the
stabiliser of its character, and write $d_\rho=p^k$. The character
eigenspaces of $\rho|_K$ are permuted transitively by $G$, by
irreducibility and normality, and only nontrivial characters of $K$
occur. Write $\bar H_n$ for the image of $H$ modulo $p^n$. Every
$\sigma\in H$ acts trivially on the $p$th roots of unity: if $\sigma$
sends a constituent $\chi$ of $\rho|_K$ to $\chi^a$, then $\chi^a=g\chi$
for some $g\in G$, and iterating gives $g^j\chi=\chi^{a^j}$; the order of
$g$ is a power of $p$, whereas $a\in\F_p^\times$ has order dividing $p-1$,
so $a=1$ (for $p=2$ this is automatic). Choose $w\in W/W^{p^n}$ with
$\chi(w^{p^{n-1}})\ne1$. It preserves the $\chi$-eigenspace, and all its
eigenvalues there are primitive $p^n$th roots of unity. This multiset is
$H$-invariant, because conjugation by $H$ preserves $\chi$ and the
isomorphism type of $\rho$. The action of $\bar H_n$ on the primitive
roots is free, so $\abs{\bar H_n}$ divides each eigenspace dimension and
hence $d_\rho$. Write $\abs{\bar H_n}=p^s$, so $s\le k$.

\emph{The orbit norm, $p$ odd.} Let $\alpha=\det D_\rho$. The contribution
of the Galois orbit of $\rho$ is $N_{L^H/\Q}(\alpha)^{d_\rho}$. By
cyclotomic local reciprocity, for $B\le A$ the unit part $u$ of a norm
from $L^B\otimes\Q_p$ lies in $B$ modulo $p^M$: the reciprocity map sends
$p$ to the identity and a unit to its residue or its inverse, and norms
from the fixed field map into $B$. Thus
$u\bmod p^n\in\bar H_n\subseteq1+p\Z/p^n\Z$. This cyclic group has a
unique subgroup of order $p^s$, so $u\equiv1\pmod{p^{n-s}}$, and
$u^{p^k}\equiv1\pmod{p^{n-s+k}}$, in particular modulo $p^n$.

\emph{The orbit norm, $p=2$, $n\ge2$.} Realness puts $\alpha$ in $L^J$
with $J=\langle H,-1\rangle$, and the orbit contribution is
$N_{L^J/\Q}(\alpha)^{2^k[J:H]}$. Let $u$ be the unit part of this norm and
write $u=\pm v$ with $v\in1+4\Z_2$. The image of $J$ contains $-1$, so
$v\bmod2^n$ lies in its principal-unit subgroup, whose order is half the
order of the image. If $J=H$, this subgroup has order $2^{s-1}$; then
$k\ge s\ge1$, the sign disappears under the exponent, and
$v^{2^k}\equiv1\pmod{2^{n+1-s+k}}$. If $J\ne H$, the subgroup has order at
most $2^s$, and the exponent $2^{k+1}$ removes the sign and gives the same
precision. These follow from the cyclic structure of
$(1+4\Z_2)/(1+2^n\Z_2)$ and the gain of one bit per squaring on $1+4\Z_2$.
For $p=2$, $n=1$ the assertion says only that the units are odd.

Multiplying the orbit contributions, and noting that $\abs K$ is a power
of $p$, proves the congruences for $U_n/U_{n-1}$ and hence for
$U_n-U_{n-1}$. Convergence and the precision of the limit follow by
summing the successive differences.
\end{proof}

\begin{proposition}[Sharpness at every level; \stU]\label[proposition]{gra:prop:unitsharp}
Let $X$ be the bouquet with $(r-1)^2$ loops of voltage $1$ and $r$ loops of
voltage $2$ in the $\Z_p$-tower, so $d=1$. Take $r=p-1$ for $p\ge5$,
$r=5$ for $p=3$, and $r=3$ for $p=2$. Then for every $n\ge1$ the ratio
$U_n/U_{n-1}$ differs from $1$ by valuation exactly $n$ for odd $p$, and by
valuation exactly $n+1$ for $p=2$ and $n\ge2$. For $p=2$, $n=1$, the
bouquet of three loops of voltage $1$ has $U_1/U_0=3$. Subdividing every
edge into a path of length $p^2$ gives towers of simple graphs with the
same $U_n$.
\end{proposition}

\begin{proof}
Using $\grah(z^2)=\grah(z)(4-\grah(z))$, the nontrivial Laplacian factor at
a character value $z$ is
\[
 (r-1)^2\grah(z)+r\grah(z^2)=\grah(z)(r+z)(r+z^{-1}).
\]
For the primitive roots of order $p^n$, whose number is even except at
order two, the product over the layer is $p^2\Phi_{p^n}(-r)^2$; the same
formula holds at order two by direct evaluation. After the matrix-tree
division by $p$ between successive levels, the unit ratio is
$\graunit(\Phi_{p^n}(-r))^2$.

For odd $p$ write $(-r)^{p^{n-1}}=1+x$; lifting the exponent gives
$v_p(x)=n$. The geometric sum yields
\[
 \frac{\Phi_{p^n}(-r)}p=1+\frac{p-1}2x
 +\sum_{j=2}^{p-1}\frac1p\binom p{j+1}x^j .
\]
Every term of the displayed sum has valuation at least $n+1$: the last
has valuation $(p-1)n-1$ and the others at least $2n$. The one exception
to this estimate is $p=3$, $n=1$, where $\Phi_3(-5)/3=7$ and
$v_3(7^2-1)=1$. In all cases $\Phi_{p^n}(-r)/p$ is a unit, and its square
differs from $1$ by valuation exactly $n$.

For $p=2$, $r=3$, $n\ge2$, the unit is $(3^{2^{n-1}}+1)/2$. Its
difference from $1$ has valuation $n$ and its sum with $1$ has valuation
$1$, so its square differs from $1$ by valuation $n+1$. For $p=2$, $n=1$,
three loops of voltage $1$ give the unit ratio $3$, which is sharp modulo
$2$ in \cref{gra:thm:units}. Since later levels are more divisible, these
examples also show that the stated precision for the limit is optimal.

Subdivide each edge, equivariantly and with separate internal vertices,
into a path of length $p^2$; a loop becomes a cycle of length at least
four. Replacing every edge by a path of length $\ell$ multiplies the
spanning-tree count by $\ell^{\abs E-\abs V+1}$, as in
\cref{gra:rem:simplebase}, so the unit parts do not change.
\end{proof}

In the abelian case, cyclic and iterated multivariable resultants govern
these counts \cite{UekiYoshizaki2022Cyclic,Yoshizaki2025Iterated}, and
the case $r=2$ of the step pattern above is the Fibonacci tower of
Kobayashi's graph class-number programme \cite{Kobayashi2025Graph}. The
argument above uses Galois stabilisers and norm groups instead of
individual cyclotomic valuations, and the representation argument
replaces the scalar characters of the abelian setting. Exact integer
Laplacian cofactors for two nonabelian covers at $n=2$ give
$(U_2,U_1)\equiv(1,10)\pmod{27}$ on $81$ vertices and
$(U_2,U_1)\equiv(13,5)\pmod{16}$ on $64$ vertices
(\artifact{spanning-tree-unit-stabilization/verification.txt}).

Determinants of Grover walks on abelian towers have valuations of the same
kind; they are computed by the automaton of \cref{gra:sec:effective}, and on
the ternary square grid at every level by \cref{gra:thm:grover} in
\cref{sec:grw}.

\subsection{Open questions}\label{gra:sec:open}
{\raggedright
\sources{Effective_Graph_Towers; grh_method_transfers_graph_towers_and_finite_rings; grh_method_transfers; grh_transfers_lubin_tate_and_graph_towers; grh_transfer_extensions_2026-11-26; five_step_graph_towers_2026-11-26; five_step_sharpness_2026-11-27; seven_step_graph_towers_2026-11-27; five_step_refinement_2026-11-27; sparse-padic-roots-and-graph-towers}
}

\begin{enumerate}
\item \emph{The sharp constant for the onset.} For a fixed prime and a
fixed loop count $m\ge\max(3,p)$, what is the sharp coefficient of the
voltage bit length $b$ in the worst-case onset $L_{p,m}(b)$ of
\cref{gra:thm:loops}? \Cref{gra:cor:logbound} gives at most
$3b\log_p2+O_p(1)$, and the constructions give at least
$b\log_p2-O_p(1)$.
\item \emph{Higher rank.} For fixed $p$, $d\ge3$ and a fixed number of
voltage vectors, what is the largest possible onset for vectors of bit
length at most $b$? In rank three the exceptional directions can meet
along a positive-dimensional set, and an effective bound would require
controlling such intersections; the only general bound is that of
\cref{gra:rem:onsetbound}.
\item \emph{Minimal linear representations.} Can one compute the smallest
rational linear representation of the shell sums directly from the
exponent vectors, without constructing all reachable digit states? The
five examples of \cref{gra:thm:five} have far fewer poles than registers,
and a structural description of the cancelled factors could yield a
practical improvement for Laplacians with many terms.
\item \emph{Input models.} Which exact input representations beyond
rational voltages support a terminating algorithm as in
\cref{gra:thm:effective}? More general algebraic $p$-adic inputs would
need additional control of carries or an alternative way of recognising
coincident exponents; \cref{gra:cor:oracle} shows that an input model
matters.
\item \emph{Extremal moment coefficients.} In rank two the maximum of
$\graA(S)$ is known for three steps at every $p>6$ (\cref{gra:thm:three}),
for four steps at every $p>8$ (\cref{gra:thm:four}), for even $m$ at the
primes with $2m\mid p-1$ (\cref{gra:cor:layers}), and for five steps at
$p=11,31$, at every prime up to $181$ and at $761$, $1031$, $1511$,
$1571$, $2281$, $2351$ (\cref{gra:thm:fivestep,gra:prop:fivelarge}).
Determine the five-step maximum at the other primes; outside $p=11,31$ it
is at most $4(p+1)+16$, and for $p\equiv1\pmod{10}$ at least $4(p+1)+12$,
with $4(p+1)+14$ attained by simple towers at infinitely many such primes.
Determine
the seven-step maximum at general $p>14$, where it is at most $6(p+1)+38$
(\cref{gra:thm:seven}), the maximum for even $m\ge6$ with $2m\nmid p-1$,
and the exact maximum for general odd $m$ and $p$.
\item \emph{Distributions.} For three and four steps the distribution of
$\graA$ over all row planes is given by
\cref{gra:thm:three,gra:thm:census} through point counts of a conic and of
the curve $C$, that is, through the traces of $E_1$ and $E_2$
(\cref{gra:thm:jacobian}); for four steps the mean and variance are known
in every rank (\cref{gra:prop:meanvar}). Determine the rank-three
four-step distribution, where each row space cuts out a plane section of
$C$ and the incidence of several special points on one plane matters.
Which root configurations of the quartic, sextic and higher directional
moments can occur on a single plane for $m\ge5$? In higher dimension, can
the distribution of $\graA$ for random step lists be expressed through
point counts of the successive moment varieties, and which higher moment
curves have equally explicit rational elliptic factors?
\end{enumerate}
